\section{Supporting contractions and finite ordinary programs}
\label{sec:perturbation}\label{sec:contractions}

The generalized pairs in this section have an effective real boundary and
a fixed nef Bott--Chern class on a smooth carrier projective over the
initial space. Their structure boundaries on higher models define
generalized log discrepancies; generalized klt means that these
discrepancies are positive. Equalities of $(1,1)$-classes mean
Bott--Chern equality. Equalities of line bundles are actual isomorphisms
after clearing the stated denominators.

The proof uses Theorem~1.1 of \cite{OAI:Termination-of-generalized-log-canonical-flips-on-compact-Kahler-fourfolds-October-5-2026} only for existing
rational flip sequences. That theorem uses the low-place, relative-program,
extraction, and special-termination results proved in
Lemma~\ref{lem:low-places} and
Sections~\ref{sec:relative}--\ref{sec:special}. Those results precede
this use logically and do not depend on Theorem~\ref{thm:finite}.
In particular, the lower-dimensional
termination theorem used below is Theorem~\ref{thm:gklt-three}, proved
from the projective relative constructions in Section~\ref{sec:relative}.
When the ordinary-step construction is used in lower dimensions, it is
applied only after the corresponding assertion $\mathsf B_d$ has been
established in the induction.

\subsection{Projectivity when a rational line bundle detects the ray}

We use the cone duality and finite-dimensional real Bott--Chern
spaces of Section~\ref{sec:analytic}. A rational line bundle means an
element of $\Pic(T)\otimes\Q$. Relative ampleness always means a
single global line bundle, after clearing a denominator.

\begin{lemma}\label{repair:lem:polarization}
Let $f:T\to Z$ be a proper contraction between normal compact
K\"ahler spaces. Suppose $\gamma$ is K\"ahler on $Z$ and
\[
 \alpha=f^*\gamma,\qquad
 \NA(T)\cap\alpha^\perp=R
\]
is a nonzero ray. If a global rational line bundle $D$ has
$D\cdot R<0$, then $-D$ is relatively ample and $f$ is projective.
\end{lemma}
\begin{proof}
Fix a K\"ahler class $\omega$ on $T$ and normalize $\NA(T)$ by
$\omega\cdot z=1$. The resulting slice $S$ is compact. Its unique
point on $R$ has negative $D$-degree, so $d=c_1(D)$ is negative on
a neighborhood of that point in $S$. On the complementary compact
set, $\alpha$ has a strictly positive minimum, while $d$ is bounded.
For all sufficiently small $s>0$, the class $\alpha-sd$ is therefore
strictly positive on $S$ and is K\"ahler by cone duality.

Choose an integer $m>0$ making $mD$ integral. Then
\[
 c_1(-mD)+f^*((m/s)\gamma)=(m/s)(\alpha-sd)
\]
is K\"ahler. On a neighborhood in the base with a potential for
$\gamma$, absorb that potential in a smooth metric on the same global
line bundle $-mD$. Its curvature is positive on the fibres. The analytic
positive-line-bundle criterion and the fibrewise relative-ampleness
criterion make $-mD$ relatively ample. The line bundle was global
throughout, so no gluing of unrelated local polarizations is required.
\end{proof}

If, as in Proposition~\ref{prop:ordinary-step}, $\alpha-c_1(D)$ is already K\"ahler,
one can apply this metric argument directly with $s=1$.
This proves projectivity after the contraction exists; it does not
construct that contraction.

\subsubsection{Constructing a generalized flip from one global line bundle}

The next argument supplies a global flip algebra whenever a projective
negative contraction has an actual line bundle detecting its ray. It will
allow the lower-dimensional reductions to use ordinary rational flips.

\begin{proposition}[A global algebra for a detected flip]
\label{repair:prop:detector-flip}
Let $(T,B+\mathbf M)$ be a gklt pair on a normal compact K\"ahler
space which is globally strongly $\Q$-factorial, and let $A\in\BC(T)$
be its adjoint class. Let $f:T\to Z$ be a projective small contraction
onto a normal compact K\"ahler space. Suppose that the classes of all
$f$-vertical curves lie on one ray $R\subset\NA(T)$, that $A\cdot R<0$,
and that a global line bundle $L$ satisfies $L\cdot R<0$.

Then $\bigoplus_{m\geq0}f_*L^m$ is locally finitely generated. The
relative Proj of a sufficiently divisible Veronese is a projective
small contraction $f^+:T^+\to Z$, where $T^+$ is compact K\"ahler and
globally strongly $\Q$-factorial. The reflexive transform $L^+$ is a
rational line bundle, with a positive power equal to the relatively
ample tautological bundle. The transformed boundary and the same
b-nef datum define a gklt pair on $T^+$, with relatively K\"ahler
adjoint $A^+$. This is the generalized flip, with the usual
strict discrepancy increase at centres in either exceptional locus.

Moreover, if $t=(A\cdot C)/(L\cdot C)>0$ for an $f$-vertical curve $C$,
then there is a unique $\eta\in\BC(Z)$ such that
\[
 A=t\,c_1(L)+f^*\eta,\qquad
 A^+=t\,c_1(L^+)+(f^+)^*\eta.
\]
Every class on $T$ has the corresponding forward transport with an
actual exceptional-divisor correction. No assertion of surjectivity
on Bott--Chern groups is required.
\end{proposition}

\begin{proof}
Since $f$ has a global relatively ample bundle and all its vertical
curves lie on $R$, numerical invariance of ampleness on the projective
fibres shows that $-L$ is $f$-ample.

\smallskip\noindent
\emph{A local rational ordinary adjoint.}
Fix $z\in Z$ and a relatively compact Stein neighbourhood $U$.
Take a projective log resolution $p:V\to T_U$ carrying $\mathbf M$,
and put $\rho=f p$. Write its actual structure boundary as $B_V$, so
\[
 [K_V+B_V]+[\mathbf M_V]=p^*A.
\]
Every coefficient of $B_V$ is less than one. Canonical divisors and a
Cartier representative $D_L$ of $L$ can be chosen over $U$: the relevant
proper direct images have rank one on the birational isomorphism locus,
and Cartan's Theorem~A supplies sections nonzero there. Choose compatible
canonical representatives and set
\[
 N=t p^*D_L-K_V-B_V.
\]
For every $\rho$-vertical curve, its degree equals that of $\mathbf M_V$.
Thus $N$ is $\rho$-nef.

It is also $\rho$-big. Indeed, write $N$ as a positive convex combination
of finitely many rational Cartier divisors $N_j$ in its finite Cartier
span. If $P$ is $\rho$-ample and $q_jN_j$ is integral, Cartan's
Theorem~A applied to the rank-one sheaf $\rho_*(q_jN_j-P)$ gives
$q_jN_j\sim P+E_j$ with $E_j\geq0$. Hence $N\sim_{\R}H+E$ with $H$
relatively ample and $E\geq0$.
Choose $\epsilon>0$ small enough that $(V,B_V+\epsilon E)$ is sub-klt.
The divisor $(1-\epsilon)N+\epsilon H$ is relatively ample. Express it
as a positive combination of rational ample divisors and take general
members of sufficiently high multiples. Relative Bertini
\cite[Theorem~2.20]{DHP2024}, applied on a log resolution of $B_V+E$,
gives an effective $\Theta\sim_{\R}N$ such that $(V,B_V+\Theta)$ is
sub-klt. The high multiples make the added coefficients arbitrarily
small. All choices are made near the compact fibre; shrink $U$ once
to retain them.

Put $\Delta=p_*(B_V+\Theta)\geq0$. Pushing down the finite
principal-divisor expression and pulling it back again give
\[
 K_{T_U}+\Delta\sim_{\R}tD_L,
 \qquad p^*(K_{T_U}+\Delta)=K_V+B_V+\Theta.
\]
Thus $(T_U,\Delta)$ is klt. Any finitely many base line bundles in a
relative linear equivalence can first be trivialized by shrinking $U$.
Record the resulting equality as
\begin{equation}\label{repair:eq:detector-rationalization}
 K_{T_U}+\sum_i d_iD_i=tD_L+\sum_j b_j\operatorname{div}(g_j),
 \qquad d_i>0,
\end{equation}
using the positive support of $\Delta$. After a further relatively
compact shrinking, the displayed divisors have finitely many prime
components: their locally finite supports meet a compact inverse image.
Equality of coefficients in \eqref{repair:eq:detector-rationalization} is a
finite rational affine system in $(d_i,t,b_j)$. The klt condition is
open within this system, as seen on one log resolution; the pullbacks
of its adjoints vary linearly by the displayed identity. A nearby
rational solution therefore gives a rational effective klt boundary
$\Delta_q$ and $t_q\in\Q_{>0}$ with
\begin{equation}\label{repair:eq:detector-local-adjoint}
 K_{T_U}+\Delta_q\sim_{\Q}t_qD_L.
\end{equation}
The individual local primes $D_i$ need not be $\Q$-Cartier: the identity
ensures that the total adjoint is $\Q$-Cartier.

\smallskip\noindent
\emph{One global flip algebra.}
The adjoint in \eqref{repair:eq:detector-local-adjoint} is $f_U$-antiample.
Its ordinary flip and local finite generation follow from
\cite[Theorems~1.14 and~1.18]{Fujino2022}. Clearing the displayed
linear equivalence identifies a Veronese of its canonical algebra
with a Veronese of $\mathcal R:=\bigoplus_{m\geq0}f_*L^m$ over $U$.
Hence $\mathcal R$ is locally finitely generated by
\cite[Lemma~2.26]{Fujino2022}. The graded pieces are coherent, and
the section algebra over a connected open set is an integral domain,
as is seen using a meromorphic generator of $L$.
Compactness supplies one sufficiently divisible Veronese
$\mathcal R^{(r)}$ generated in degree one. Its relative Proj $T^+$
restricts to the ordinary flip over every such $U$. It is therefore
normal and small over $Z$, and has the global relatively ample bundle
$\mathcal O_{T^+}(1)$. On the common big open this bundle is $L^r$.
Reflexivity gives
\[
 (L^+)^{[r]}\simeq\mathcal O_{T^+}(1).
\]
This proves that $L^+$ is a rational line bundle before any assertion
of factoriality. Relative positivity over the compact K\"ahler base
makes $T^+$ compact K\"ahler.

For any global rank-one reflexive sheaf $\mathcal F^+$ on $T^+$,
take its coherent reflexive transform $\mathcal F$ on $T$ using a
common resolution. Some $M=\mathcal F^{[m]}$ is a line bundle.
Choose $c\in\Q$ so that $(M+cL)\cdot C=0$, and clear its denominator
to obtain an integral bundle $J=n(M+cL)$ numerically trivial over $f$.
The local rational klt pairs \eqref{repair:eq:detector-local-adjoint} satisfy
$J-(K_{T_U}+\Delta_q)$ relatively nef. Integral descent (Lemma~\ref{lem:descent}) yields
one global line bundle $J_Z$ with $J=f^*J_Z$.
On $T^+$, pull back $J_Z$, undo the rational twist $ncL^+$, and compare
on the common big open. After clearing the already established index
of $L^+$, reflexivity gives an invertible positive reflexive power of
$\mathcal F^+$. This proves global strong factoriality.

\smallskip\noindent
\emph{The original generalized pair.}
The local log-Fano pairs above give $R^if_*\mathcal O_T=0$ for $i>0$
by relative vanishing. They make $T$ locally klt, hence rational.
Leray for a projective resolution then shows that $Z$ is rational.
By \cite[Lemma~8.7]{DHP2024}, $f^*\BC(Z)=R^\perp$ and pullback is
injective. Thus $A=t c_1(L)+f^*\eta$ for a unique $\eta$.
Define $A^+=t c_1(L^+)+(f^+)^*\eta$; it is relatively K\"ahler.

To construct its actual discrepancy divisor, let $\mathcal I$ be the
ideal defined by
\[
 \operatorname{im}(f^*f_*L^r\longrightarrow L^r)
   =\mathcal I\otimes L^r.
\]
Principalize $\mathcal I$ on a common resolution $p:W\to T$,
$q:W\to T^+$ carrying $\mathbf M$. If
$\mathcal I\mathcal O_W=\mathcal O_W(-F_L)$, the tautological quotient
gives the actual bundle identity
\[
 q^*\mathcal O_{T^+}(1)=p^*L^r\otimes\mathcal O_W(-F_L).
\]
Hence $E_L=F_L/r$ is effective and exceptional over both sides, and
$[E_L]=p^*c_1(L)-q^*c_1(L^+)$. If $B_W$ is the original structure
boundary, put $B_W^+=B_W-tE_L$. Then
\[
 [K_W+B_W^+]+[\mathbf M_W]=q^*A^+.
\]
Its pushforward is $B^+$, and its discrepancies do not decrease.
The local ordinary comparison in \eqref{repair:eq:detector-local-adjoint}
has discrepancy divisor $t_qE_L$, so its strictness proves the stated
strictness for $tE_L$ as well. Thus the original pair remains gklt.

Finally, write any $\theta\in\BC(T)$ uniquely as
$\theta=f^*\eta_\theta+s c_1(L)$ and set
$\theta^+=(f^+)^*\eta_\theta+s c_1(L^+)$. Its correction is the actual
divisor $sE_L$. For a global divisor class this is its strict transform:
the difference between the two candidate corrections is exceptional
and numerically trivial over $T^+$, hence zero by negativity.
\end{proof}

\subsection{Ordinary real boundaries and forward transport}\label{repair:sec:real}

Its rational-step input is the proof of Proposition~\ref{prop:ordinary-step},
with its supporting contraction supplied by
Proposition~\ref{repair:prop:supporting-contraction} below and its projectivity
supplied by Lemma~\ref{repair:lem:polarization}.

\begin{proposition}\label{repair:prop:real}
Let $T$ be normal compact K\"ahler and globally Weil $\Q$-factorial,
with canonical sheaf a rational line bundle. Let $B\geq0$ be a real
boundary with $(T,B)$ klt. Assume the rational ordinary-step result
of Proposition~\ref{prop:ordinary-step} for $T$ and its subsequent models. Then every
$(K_T+B)$-negative extremal analytic ray has an ordinary step with
projective contractions, compact K\"ahler models, and the expected
opposite relative ample signs. Global Weil factoriality is preserved;
global strong factoriality is preserved when imposed initially.
\end{proposition}
\begin{proof}
Put $D=K_T+B$ and fix a $D$-negative extremal ray $R$. On the finite
positive support of $B$, effectiveness, the klt condition and negativity
on $R$ persist under small coefficient changes. Klt openness is checked
on one log resolution of this support. Choose nearby effective rational
klt boundaries $B_0,\ldots,B_m$ and positive real numbers $u_j$ with
\[
 B=\sum_j u_jB_j,\qquad \sum_j u_j=1,
 \qquad D_j\cdot R<0,\quad D_j=K_T+B_j.
\]
For $B=0$ take just $B_0=0$. Every $D_j$ is a global rational line
bundle. The rational-step theorem for $(T,B_0)$ provides a projective
contraction $f:T\to Z$, a compact K\"ahler base, and
\[
 f^*\BC(Z)=R^\perp,\qquad \rho(T/Z)=1,
 \qquad -D_0\text{ relatively ample}.
\]
Let $C$ be a rational curve spanning $R$, and set
$a_j=(D_j\cdot C)/(D_0\cdot C)\in\Q_{>0}$. On every projective
fibre, $-D_j$ is numerically equivalent to the positive multiple
$-a_jD_0$. Numerical invariance of ampleness on that fibre, followed
by the relative-ampleness criterion, makes $-D_j$ relatively ample.
The positive combination $-D$ is relatively ample as a real line
bundle. This also handles a fibre-type contraction.

Suppose henceforth that $f$ is birational. A Cartier multiple
$M_j=n_j(D_j-a_jD_0)$ is numerically trivial over $Z$ and
$M_j-D_0$ is relatively nef. Integral descent, Lemma~\ref{lem:descent},
applied to the rational klt pair $(T,B_0)$ gives a rational line bundle
$A_j$ on $Z$ with the actual identity
\begin{equation}\label{repair:eq:descent}
 D_j=a_jD_0+f^*A_j.
\end{equation}
Here and below an identity of rational line bundles means an isomorphism
after a common integral multiple. The descent lemma is used only for
birational $f$.

For a small $f$, let $f^+:T^+\to Z$ be its rational $D_0$-flip.
The rational theorem supplies a global relatively ample adjoint
$D_0^+$ and preserves the stated factoriality and canonical-sheaf
conditions. The other transformed adjoints $D_j^+$ are therefore
rational line bundles. Transform \eqref{repair:eq:descent} on the common
big open and extend by reflexivity to obtain
\[
 D_j^+=a_jD_0^++(f^+)^*A_j.
\]
All $D_j^+$ are relatively ample, as is
$D^+=\sum_j u_jD_j^+$. Thus this same small map is the required
ordinary real-boundary flip. The real ordinary discrepancy comparison
proves that $(T^+,B^+)$ is klt and gives the strict increases at
exceptional centres.

For a divisorial contraction, write
$D_0=f^*D_0'+e_0E$ with $e_0>0$. Transforming
\eqref{repair:eq:descent} on the target and summing yields
\[
 D-f^*D'=\left(\sum_j u_ja_j\right)e_0E.
\]
This coefficient is positive. The same discrepancy comparison proves
the required assertions. The ambient category was already preserved
by the rational step in both cases.
\end{proof}

\subsubsection{Forward transport only}
For a birational step of Proposition~\ref{repair:prop:real}, write
$f':T'\to Z$ for the positive contraction, or the identity in the
divisorial case. Negative-side rank one gives, uniquely,
\[
 \theta=f^*\eta+s\,c_1(D),\qquad
 \theta'=(f')^*\eta+s\,c_1(D').
\]
On a common resolution $p:V\to T$, $q:V\to T'$,
\[
 p^*\theta-q^*\theta'=s[p^*D-q^*D'].
\]
The expression in brackets is the actual exceptional discrepancy
divisor. This proves the forward transport used in Lemma~\ref{lem:class-transport}. If $p^*H_T=p_0^*H+[J_T]$ is the fixed-nef-datum relation,
then
\[
 J_{T'}=J_T-s(p^*D-q^*D'),\qquad
 q^*H_{T'}=p_0^*H+[J_{T'}].
\]
No divisor is extracted, so $J_{T'}$ is $q$-exceptional. Its negative
is $q$-nef because $p_0^*H$ is nef; negativity gives $J_{T'}\geq0$.
There is no assertion of a surjection onto all of $\BC(T')$ or of
positive-side relative rank one.

\subsection{Ingredients for the contraction theorem}
\label{repair:subsec:bounded-bpf-induction}

We isolate the part of the transcendental base-point-free argument that
is needed here. The induction below uses ordinary rational termination
only in dimensions at most three. In particular, it does not use a
finiteness theorem for generalized minimal models.

For an integer $d\geq0$, consider the following statements, each in
dimensions at most $d$.
\begin{description}
\item[$\mathsf B_d$.]
If $(X,B+\mathbf M)$ is a compact K\"ahler gklt pair,
$B+\mathbf M_X$ is modified big, and
$\alpha=[K_X+B+\mathbf M_X]$ is nef, then there is a Moishezon
contraction $h:X\to Y$ onto a compact K\"ahler space and a K\"ahler
class $\omega_Y$ such that $\alpha=h^*\omega_Y$.
\item[$\mathsf C_d$.]
Let $(X,B+\mathbf M)$ be a compact K\"ahler generalized pair, with
$B+\mathbf M_X$ modified big and
$\alpha=[K_X+B+\mathbf M_X]$ nef and NQC. If the restriction of
$\alpha$ to the closed subspace defined by
$\mathcal J(X,B+\mathbf M)$ has a Moishezon contraction, then
$\alpha$ has a Moishezon contraction.
Write $\mathsf C_{d,\mathrm{big}}$ for this assertion with the
additional assumption that $\alpha$ is big.
\item[$\mathsf G_d$.]
If $(X,B+\mathbf M)$ is a globally strongly $\Q$-factorial compact
K\"ahler gklt pair, $B+\mathbf M_X$ is modified big, and its adjoint
class $A$ is pseudo-effective, then it has a good minimal model.
One may obtain this model from a smooth modification by finitely many
ordinary rational klt steps. The induced map from $X$ extracts no
divisor and is $A$-nonpositive.
\end{description}
Here ``NQC'' has the cohomological meaning of
\cite[Definition~2.43]{HX2026}: the class is a positive real
combination of rational classes in $H^2$ that are nonnegative on
curves. All uses below require only this weak version. A Moishezon
contraction has the meaning of \cite[Definition~2.32]{HX2026}; in
particular its fibres are connected by chains of compact curves.

We first record a polarization fact for actual line bundles. It is
the only promotion from a Moishezon map to a projective map used in
the non-klt gluing construction.

\begin{lemma}[Polarization by an actual real line bundle]
\label{repair:lem:actual-line-polarization}
Let $f:T\to Z$ be a proper Moishezon map, where $T$ is a compact
K\"ahler space. Suppose that $L\in\operatorname{Pic}(T)\otimes\R$
and a K\"ahler class $\omega$ satisfy
\[
 L\cdot C=\omega\cdot C
 \quad\text{for every compact curve $C$ contracted by $f$.}
\]
Then $L$ is relatively ample and $f$ is projective.
\end{lemma}

\begin{proof}
Let $V$ be an irreducible positive-dimensional subspace of a reduced
fibre. It is Moishezon. Normalize $V$ and take a smooth projective
modification $\pi:\widetilde V\to V^{\mathrm n}$, chosen so that
an effective exceptional divisor $E$ has $-E$ relatively ample.
The pullback of the K\"ahler class from $V$ to its normalization is
K\"ahler. Consequently
$\pi^*\omega-\delta[E]$ is K\"ahler for sufficiently small
$\delta>0$. On the projective manifold $\widetilde V$ the actual
real line bundle $\pi^*L-\delta E$ has the same curve degrees as
that K\"ahler class, and is therefore ample. For example, subtract
a sufficiently small positive multiple of a fixed ample class from
the K\"ahler class; the corresponding real line bundle is nef by
the projective numerical criterion. Also $\pi^*L$ is nef, and the
decomposition
\[
 \pi^*L=(\pi^*L-\delta E)+\delta E
\]
shows that it is big. Thus
\[
 L^{\dim V}\cdot V=(\pi^*L)^{\dim V}>0.
\]
The real Nakai--Moishezon criterion for proper algebraic spaces
\cite[Theorem~1.6]{FujinoMiyamoto2020}, applied to the algebraizations
of the Moishezon fibres, proves that $L$ is ample on every reduced
fibre. The criterion is insensitive to nilpotents and reducible
components. Fibrewise ampleness is open for a proper analytic map;
in the finite-dimensional span of the finitely many line bundles
occurring in $L$, compactness therefore gives a neighbourhood of
$L$ whose rational points are relatively ample. This proves relative
ampleness of $L$ and supplies a global relatively ample rational
line bundle. Clearing its denominator proves projectivity.
\end{proof}

\begin{remark}
The proof uses only positivity of the top self-intersections of
$L$. Equality of curve degrees with $\omega$ does not assert
equality of their higher intersection numbers or of their
Bott--Chern classes.
\end{remark}



\subsubsection{Relative algebraicity over the fixed smooth base}
\label{repair:sec:fixed-smooth-base-algebraicity}

\begin{lemma}[A rational Hodge correction over a smooth base]
\label{repair:lem:fixed-smooth-base-algebraicity}
Let $f:T\to S$ be an already projective surjective morphism of normal
compact K\"ahler spaces. Assume that $S$ is smooth, $T$ has strongly
$\mathbb{Q}$-factorial klt singularities, and, for a projective resolution
$r:W\to T$, the morphism $g=f\circ r$ satisfies
\[
 g^*:H^0(S,\Omega_S^2)\xrightarrow{\ \simeq\ }
                         H^0(W,\Omega_W^2).
\]
Then every $\alpha\in H^{1,1}_{\mathrm{BC}}(T)$ can be written
\[
                         \alpha=c_1(L)+f^*\gamma,
\]
where $L\in\operatorname{Pic}(T)\otimes_{\mathbb Z}\mathbb R$ is a
finite real combination of global line bundles and
$\gamma\in H^{1,1}(S,\mathbb R)$.
\end{lemma}

\begin{proof}
The composition $g$ is projective, since all the spaces are compact.
Choose a $g$-ample line bundle on $W$, with integral Chern class $H$.
Set $d=\dim W-\dim S$ and
\[
                            c=g_*(H^d)>0.
\]
Thus $c$ is the positive degree of $H^d$ on a general fibre. The
cohomological pushforward here is the usual pushforward between the
smooth compact manifolds $W$ and $S$, defined over $\mathbb Q$ and of
Hodge bidegree $(-d,-d)$.

Define a rational linear operator on degree-two cohomology by
\[
 \Pi(\eta)=\eta-g^*\!\left(\frac{g_*(\eta\smile H^d)}{c}\right).
\]
Every class of type $(2,0)$ on $W$ is pulled back from $S$, and the
same holds for type $(0,2)$. The projection formula therefore shows
that $\Pi$ kills both types. It preserves type $(1,1)$. Consequently
\[
        \Pi(H^2(W,\mathbb Q))\subset
                         H^2(W,\mathbb Q)\cap H^{1,1}(W).
\]
By the Lefschetz $(1,1)$ theorem the image consists of rational Chern
classes of global line bundles. Applying $\Pi$ to $r^*\alpha$ gives
\[
 r^*\alpha=c_1(L_W)+g^*\gamma,
 \qquad L_W\in\operatorname{Pic}(W)\otimes\mathbb R,
 \qquad
 \gamma=\frac{g_*(r^*\alpha\smile H^d)}{c}\in H^{1,1}(S,\mathbb R).
\]
For example, express $r^*\alpha$ in a rational basis of $H^2(W,\mathbb Q)$
and apply $\Pi$ to each basis vector. This gives the required finite
real combination $L_W$.

Write $L_W=\sum_j a_jL_j$ with actual line bundles $L_j$ on $W$.
For each $j$, let
\[
                   \mathcal F_j=(r_*L_j)^{**}.
\]
It is a rank-one reflexive coherent sheaf. Strong
$\mathbb Q$-factoriality gives an integer $m_j>0$ for which
$M_j=\mathcal F_j^{[m_j]}$ is a line bundle on $T$.
The canonical identification over the isomorphism locus of $r$ gives
\[
             L_j^{\otimes m_j}\simeq r^*M_j\otimes\mathcal O_W(E_j)
\]
for an integral $r$-exceptional divisor $E_j$.
One may obtain this comparison directly from the evaluation map for
$r_*L_j$: the two coherent rank-one sheaves agree away from the
exceptional locus, so their invertible transforms differ by an
exceptional divisor. No global meromorphic frame of $L_j$ or of the
canonical sheaf is required.

Set $L=\sum_j(a_j/m_j)M_j$ and
$E=\sum_j(a_j/m_j)E_j$. The preceding identities yield
\[
                r^*(\alpha-c_1(L)-f^*\gamma)=[E].
\]
The divisor $E$ is $r$-exceptional and is numerically trivial on every
$r$-contracted curve. The exceptional negativity lemma applied in both
signs gives $E=0$. Injectivity of pullback by a resolution, or pushforward
of the equality of currents, now gives
$\alpha=c_1(L)+f^*\gamma$.
\end{proof}

\paragraph{What this proves about the relative curve spaces.}
On real combinations of curves contracted by $f$, numerical equivalence
tested by global line bundles is exactly numerical equivalence tested
by all Bott--Chern classes. Indeed, the displayed decomposition makes
every Bott--Chern pairing on such curves a pairing with $c_1(L)$, while
the pullback term pairs to zero. The converse implication is immediate.
The resulting finite-dimensional isomorphism identifies the closures
of the cones generated by the contracted curves, and therefore their
extremal rays. This statement does not assert that every class on an
arbitrary birational model transfers backwards, or that every
numerically trivial class descends under an arbitrary morphism.

\paragraph{Persistence during the initial relative program.}
Suppose a finite prefix of a projective relative MMP over this fixed
smooth $S$ has produced $T_i\to S$. Its morphism to $S$ is projective by
the projective relative MMP construction, and its strong
$\mathbb Q$-factorial klt property is preserved. Take a common
resolution of $T_i$ and the initial $T$ over $S$. Holomorphic two-forms
on smooth compact K\"ahler manifolds are invariant under modifications.
The isomorphism of holomorphic two-forms in the lemma therefore holds
for every projective resolution of $T_i$. The lemma applies anew on
each $T_i$; no projectivity-descent theorem is being used to establish
this persistence.

\paragraph{Application to the smooth MRC base.}
For a smooth holomorphic model $W\to S$ of the MRC fibration, the
pullback of holomorphic two-forms is an isomorphism, since the general
fibres are rationally connected. Projectivity of this smooth model is
the smooth-base, holomorphic-two-form part of the projectivity
criterion. Its rational Hodge correction proof is independent of the
singular projectivity-descent assertion. Once that projective model
has been chosen, the lemma above gives the relative real line bundles
needed for the initial MMP over $S$.

\paragraph{Local divisor representatives.}
The global objects furnished by the lemma are real combinations of
line bundles. Over each sufficiently small Stein open subset of $S$
they have actual meromorphic divisor representatives: after twisting a
line bundle by a sufficiently high power of a relatively ample line
bundle, relative generation and Cartan's theorem~A provide nonzero
sections; taking the quotient of two such sections gives a meromorphic
section of the original line bundle. Hence the ordinary projective
theorems can be used on the finite Stein cover with the intrinsic
global line bundles and their section algebras. A global effective
divisor representative is not an additional hypothesis.

\paragraph{NQC transport is a separate one-way assertion.}
\label{repair:par:one-way-nqc}
The rational Hodge correction above is not a substitute for descent
along a contracted ray. For a nef class $\alpha$ whose contraction is
$\alpha$-trivial, use its already established Bott--Chern descent
$\alpha=\pi^*\alpha_Z$. If $V_{\mathbb Q}\subset H^2(T,\mathbb Q)$ is
the minimal rational subspace containing $\alpha$, then
\[
 V_{\mathbb Q}\subset\pi^*H^2(Z,\mathbb Q),
\]
because the right-hand side is a rational subspace containing
$\alpha$. Choose the NQC summands in $V_{\mathbb Q}$: intersect the
rational polyhedral cone generated by any initial NQC summands with
$V_{\mathbb Q}\otimes\mathbb R$, and express $\alpha$ using rational
generators of this intersection. Those generators remain nonnegative
on curves. These rational NQC summands therefore descend and can be
pulled forward to the flipped model. For nonnegativity, lift a curve
on the new model to a curve on a common projective resolution that
maps onto it with some positive degree. Equality of the pulled-back
classes shows that its pairing is nonnegative; dividing by this
positive degree preserves the sign. This proves weak NQC on the new
model. It does not preserve a specified denominator, and no
degree-one lift is asserted. Below, the initial relative program
preserves its denominators by integral line-bundle descent, whereas
the lower-dimensional absolute program chooses a new bound on each
model. Neither argument requires surjectivity of the forward
Bott--Chern map or equality of the two models' Bott--Chern dimensions.

\subsubsection{Fixed relative degree denominators by integral descent}
\label{repair:subsec:relative-nqc-integral}

The initial program over the fixed smooth MRC base has a stronger
property than is needed for arbitrary birational maps: its rational
cohomology classes can be represented by rational line bundles modulo
that base. Integral descent of those line bundles preserves their
degree denominators. This avoids any assertion that a curve has a
degree-one lift through a resolution.

\begin{lemma}[Integral descent at a projective log-Fano contraction]
\label{repair:lem:relative-nqc-integral-descent}
Let $h:T\to Y$ be a projective bimeromorphic contraction of normal
complex spaces. Suppose that each point of $Y$ has a neighbourhood
$U$ on which an effective real boundary $\Theta_U$ gives an ordinary
klt pair $(h^{-1}U,\Theta_U)$ and
$-(K_{h^{-1}U}+\Theta_U)$ is $h$-ample. If $L$ is an actual line
bundle with $L\cdot C=0$ for every $h$-contracted curve, then
$h_*L$ is a line bundle and the natural evaluation map is an
isomorphism
\[
                         h^*h_*L\simeq L.
\]
In particular no additional tensor power is needed.
\end{lemma}

\begin{proof}
Fix $y\in Y$ and shrink to a Stein neighbourhood on which the boundary
is defined. A Cartier-divisor representative of $L$ exists there:
$h_*L$ is coherent and has rank one on the bimeromorphic isomorphism
locus, so Cartan's theorem~A supplies a section nonzero on that locus.
Its pullback is a nonzero meromorphic frame of $L$ and gives the
required Cartier divisor.

The line bundle $L$ is $h$-nef. Numerical invariance of ampleness on
the projective fibres gives
\[
                 L-(K_T+\Theta_U)\quad\text{$h$-ample}.
\]
The analytic base-point-free theorem
\cite[Theorem~6.2]{Fujino2022} applies to this Cartier divisor and the
real klt boundary. It states that $L^m$ is relatively generated for
\emph{every} sufficiently large integer $m$ near the compact fibre
$h^{-1}(y)$, rather than only for divisible $m$.

Choose finitely many generating sections. The resulting map from
the reduced projective fibre to projective space is constant on each
irreducible component, since a positive-dimensional image would give
a curve of positive $L^m$-degree. The fibre is connected, so these
constant values agree. A suitable linear combination of the sections
therefore has no zero on the reduced fibre. In the local ring at a
point of the possibly nonreduced fibre it is a unit, so it is a
generator there as well. Properness allows the base neighbourhood to
be shrunk so that this section trivializes $L^m$ throughout its
inverse image.

Apply this to two consecutive sufficiently large integers $m,m+1$.
Taking the quotient of the two trivializations trivializes $L$.
Since $h_*\mathcal O_T=\mathcal O_Y$, it follows locally that
$h_*L\simeq\mathcal O_Y$ and evaluation is an isomorphism. These are
intrinsic assertions and hence glue over $Y$. This also explains why
possible torsion in a fibre Picard group introduces no new index.
The conclusion is recorded directly for extremal contractions in
\cite[Theorem~7.2(2)(iii)]{Fujino2022}; its proof allows real
boundaries. The argument above gives the intrinsic global form needed
here and does not require the contraction to have relative Picard
rank one in every local analytic neighbourhood.
\end{proof}

\begin{proposition}[A uniform relative NQC bound]
\label{repair:prop:relative-nqc-uniform}
Let $f_0:U\to Z$ be an already projective surjective morphism with
connected fibres between smooth compact
K\"ahler manifolds, and suppose that
\[
 f_0^*:H^0(Z,\Omega_Z^2)\longrightarrow H^0(U,\Omega_U^2)
\]
is an isomorphism. Let $\alpha_0$ be a nef class with an expression
\[
       \alpha_0=\sum_{j=1}^k r_j\eta_{j,0},\qquad
       r_j>0,\quad \eta_{j,0}\in H^2(U,\mathbb Q),
\]
where $\eta_{j,0}$ has nonnegative degree on every curve vertical
over $Z$. Let $D_0=K_U+\Delta_U+\mathbf M_U$ be a generalized klt
adjoint. Consider its projective relative program for
$D_0+a\alpha_0$, under the projective local ordinary reduction described in the
proof of Proposition~\ref{repair:prop:restricted-bpf-induction}.

If $\alpha_0=0$, triviality is immediate. Otherwise there are
positive integers $m_j$, chosen once on $U$, and a number
\[
                         \delta=\min_j\frac{r_j}{m_j}>0
\]
such that, for $a\delta>2\dim U$, the entire relative program is
$\alpha_0$-trivial. The same number $a$ works at every step.
\end{proposition}

\begin{proof}
\emph{Rational classes modulo the fixed base.}
Choose an $f_0$-ample line bundle with integral Chern class $H$, and
put $e=\dim U-\dim Z$ and $c=f_{0*}(H^e)>0$. The rational operator
\[
 \Pi(\eta)=\eta-
 f_0^*\!\left(\frac{f_{0*}(\eta\smile H^e)}{c}\right)
\]
preserves type $(1,1)$ and kills types $(2,0)$ and $(0,2)$.
Indeed these latter classes are pulled back from $Z$, and the
projection formula applies. Thus
$\Pi(H^2(U,\mathbb Q))\subset H^{1,1}(U)\cap H^2(U,\mathbb Q)$.
By the Lefschetz $(1,1)$ theorem, choose an actual line bundle
$L_{j,0}$ and an integer $m_j>0$ with
\[
 \eta_{j,0}=\frac1{m_j}c_1(L_{j,0})+f_0^*\gamma_j,
 \qquad \gamma_j\in H^2(Z,\mathbb Q).
\]
The original summands $\eta_{j,0}$ need not have type $(1,1)$.
Their base components account for this. Their weighted sum gives
\begin{equation}\label{repair:eq:relative-nqc-line-expression}
 \alpha_0=\sum_j\frac{r_j}{m_j}c_1(L_{j,0})+f_0^*\Gamma,
 \qquad \Gamma=\sum_jr_j\gamma_j\in H^{1,1}(Z,\mathbb R).
\end{equation}
The final assertion follows because $f_0^*$ is an injective Hodge
map and the other terms have type $(1,1)$. The identity is therefore
also one of Bott--Chern classes.
Each $L_{j,0}$ has nonnegative integral degree on every
$Z$-vertical curve.

\smallskip\noindent
\emph{Induction across an actual contraction.}
Suppose at a finite stage $f_i:U_i\to Z$ we have actual line bundles
$L_{j,i}$, nef over $Z$, and
\[
 \alpha_i=\sum_j\frac{r_j}{m_j}c_1(L_{j,i})+f_i^*\Gamma.
\]
Assume also that $\alpha_i$ is nef and that the transformed
$D_i$ is generalized klt. These assertions hold initially.
Let $R$ be a $(D_i+a\alpha_i)$-negative extremal ray over $Z$.
It is $D_i$-negative because $\alpha_i$ is nef. The projective
relative cone theorem supplies a rational generator $C$ with
\[
                     0<-D_i\cdot C\leq2\dim U.
\]
All $L_{j,i}\cdot C$ are nonnegative integers. If
$\alpha_i\cdot C>0$, one of them is at least one, and hence
$\alpha_i\cdot C\geq\delta$. Therefore
\[
 (D_i+a\alpha_i)\cdot C
             \geq-2\dim U+a\delta>0,
\]
a contradiction. We have $\alpha_i\cdot C=0$ and
$L_{j,i}\cdot C=0$ for every $j$. Every curve contracted by the
projective extremal contraction $h_i:U_i\to Y_i$ has class on
$R$ when tested by global line bundles, so the same vanishing
holds for all of them.

On each fixed Stein chart of $Z$, choose the effective real
ordinary representative of the initial relatively ample nef
datum once, before running the program. Its actual real linear
equivalence to the fixed real line-bundle representative persists
under pushforward. The Bott--Chern comparison with the generalized
adjoint modulo the fixed base class persists separately by
exceptional negativity, as detailed below.
Thus every $h_i$ is, locally on $Y_i$, a contraction with an
ordinary real klt log-Fano boundary. This verifies exactly the
hypothesis of Lemma~\ref{repair:lem:relative-nqc-integral-descent}.
Consequently
\[
 M_{j,i}:=(h_i)_*L_{j,i}\ \text{is a line bundle},\qquad
 h_i^*M_{j,i}=L_{j,i}.
\]
For a divisorial contraction set $L_{j,i+1}=M_{j,i}$.
For a flip $h_i^+:U_{i+1}\to Y_i$, set
$L_{j,i+1}=(h_i^+)^*M_{j,i}$.
These are actual line bundles with the \emph{same} integers $m_j$.
The descended Bott--Chern class
\[
 \alpha_{Y_i}
     =\sum_j\frac{r_j}{m_j}c_1(M_{j,i})+f_{Y_i}^*\Gamma
\]
satisfies $h_i^*\alpha_{Y_i}=\alpha_i$.
It is nef by descent of nefness through the projective surjection
$h_i$; its pullback to the next model is $\alpha_{i+1}$.
This proves that the step is crepant for $\alpha_i$. Hence it is
also a $D_i$-negative step, and the original generalized pair
remains gklt.

Finally each $M_{j,i}$ is nef over $Z$. For a $Z$-vertical curve
$B\subset Y_i$, projectivity supplies a curve $B'\subset U_i$
mapping onto it with some positive degree $d_B$. Then
\[
 d_B(M_{j,i}\cdot B)=L_{j,i}\cdot B'\geq0.
\]
No assertion $d_B=1$ is required. Pullback preserves this relative
nefness, so the induction applies on $U_{i+1}$. Its curve degrees
are again integers because the transported objects are line
bundles, not because any curves were lifted with degree one.
This proves the fixed bound through the whole program.
\end{proof}

\begin{remark}[Scope of the degree grid]
The transported rational classes
\[
 \eta_{j,i}=\frac1{m_j}c_1(L_{j,i})+f_i^*\gamma_j
\]
have degrees in $m_j^{-1}\mathbb Z_{\geq0}$ on curves vertical over
the fixed smooth base $Z$. This relative assertion is exactly what
the initial projective program requires. It does not assert
nonnegativity of these individual classes on all curves of every
birational model, and uses neither rational connectedness of
resolution fibres nor Graber--Harris--Starr.
\end{remark}


\subsubsection{Adaptive bounds for the lower-dimensional program}

\begin{lemma}[Adaptive NQC bounds]
\label{repair:lem:adaptive-nqc-program}
Let $1\leq d\leq3$, and assume $\mathsf B_d$ and ordinary rational
klt termination in dimension at most $d$. Let $X_0$ be a globally
strongly $\mathbb Q$-factorial compact K\"ahler klt space of dimension
$d$, and let $\alpha_0$ be a nef, weakly NQC Bott--Chern class which
is not big. Assume that
\[
                         Q_0:=\alpha_0-c_1(K_{X_0})
\]
is big. Then there is a finite sequence of ordinary $K$-negative
divisorial contractions and flips, all $\alpha_0$-trivial, ending in
an ordinary Mori fibre space $f:X_m\to Z$. The class $\alpha_m$
descends to a nef weakly NQC class on $Z$.
\end{lemma}

\begin{proof}
We first explain the one-way NQC transport that will be used. Let
$f_i:X_i\to Z_i$ be a projective ordinary $K_{X_i}$-negative ray
contraction, and suppose that $\alpha_i$ vanishes on its ray. The
ordinary contraction theorem gives
\[
                 \alpha_i=f_i^*\beta_i.
\]
This is Bott--Chern descent on the already projective log Fano
contraction, by Lemma~\ref{lem:pullback-image}; the target has rational
singularities and the required higher direct images vanish. We will also use
injectivity of
$f_i^*:H^2(Z_i,\mathbb Q)\to H^2(X_i,\mathbb Q)$.
For clarity, it follows from relative vanishing and the usual
pluriharmonic-function sequence. Indeed,
$(f_i)_*\mathcal O_{X_i}=\mathcal O_{Z_i}$,
$R^1(f_i)_*\mathcal O_{X_i}=0$, and
$(f_i)_*\mathcal H_{X_i}=\mathcal H_{Z_i}$ imply
$R^1(f_i)_*\mathbb R=0$; the Leray sequence gives injectivity on
$H^2(-,\mathbb R)$, and hence on rational cohomology.

Choose a weak NQC expression
\[
 \alpha_i=\sum_{j=1}^{r_i}a_{ij}\xi_{ij},
 \qquad a_{ij}>0,\qquad
 \xi_{ij}\in H^2(X_i,\mathbb Q),\qquad
 \xi_{ij}\cdot C\geq0
\]
for every compact curve $C\subset X_i$, with the $\xi_{ij}$ in
the minimal rational subspace containing $\alpha_i$.
Since the image of $f_i^*$ on rational cohomology is a rational
subspace and its realification contains $\alpha_i$, that minimal
subspace lies in the image. Thus
\[
 \xi_{ij}=f_i^*\eta_{ij},\qquad
 \eta_{ij}\in H^2(Z_i,\mathbb Q),\qquad
 \beta_i=\sum_j a_{ij}\eta_{ij}.
\]
The last equality follows from injectivity. In the birational case,
write $g_i:X_{i+1}\to Z_i$ for the positive-side map, using the
identity map for a divisorial contraction, and set
\[
 \alpha_{i+1}=g_i^*\beta_i,\qquad
 \xi_{ij}^+=g_i^*\eta_{ij}.
\]
On a common projective resolution $p:W\to X_i$, $q:W\to X_{i+1}$,
we have $p^*\xi_{ij}=q^*\xi_{ij}^+$. For a curve
$C^+\subset X_{i+1}$, projectivity of $q$ supplies a curve
$\widetilde C\subset W$ mapping onto it with some degree $e>0$.
Consequently
\[
 e\,\xi_{ij}^+\cdot C^+
   =\xi_{ij}\cdot p_*\widetilde C\geq0.
\]
Thus $\alpha_{i+1}$ is weakly NQC. This uses a positive-degree
multisection, not a degree-one lift. Its rational components may
acquire different integral denominators on $X_{i+1}$.
Nefness follows from equality of the pullbacks and descent of nef
Bott--Chern classes under a proper surjective morphism. In the
fibre-type case, lifting curves from $Z_i$ through the projective
morphism $f_i$ proves weak NQC for $\beta_i$; its nefness again
follows by descent from $f_i^*\beta_i=\alpha_i$.

We now choose each step separately. At a current model $X_i$, take
any weak NQC expression as above, and choose integers $l_{ij}>0$
such that $l_{ij}\xi_{ij}$ is integral. Choose
\begin{equation}\label{repair:eq:adaptive-nqc-bound}
                    t_i>2d\max_j\frac{l_{ij}}{a_{ij}}.
\end{equation}
If $\alpha_i=0$, take any $t_i>0$ and omit the bound.
As long as $Q_i:=\alpha_i-c_1(K_{X_i})$ is big and $\alpha_i$
is not big, the class $c_1(K_{X_i})+t_i\alpha_i$ cannot be
pseudoeffective. Otherwise
\[
 (t_i+1)\alpha_i
       =(c_1(K_{X_i})+t_i\alpha_i)+Q_i
\]
would be big. The nef datum $t_i\alpha_i$ descends to $X_i$, so
adding it to the ordinary klt pair $(X_i,0)$ changes no discrepancy.
In particular this generalized klt adjoint is not nef, so the
analytic cone theorem provides a negative analytic
extremal ray $R_i$. Nefness of $\alpha_i$ makes $R_i$ an ordinary
$K_{X_i}$-negative ray. Choose its ordinary length-bounded rational
curve $C_i$, so that
\[
                       0<-K_{X_i}\cdot C_i\leq2d.
\]
If $\alpha_i\cdot C_i>0$, some $\xi_{ij}\cdot C_i$ is positive
and hence is at least $1/l_{ij}$. Then
\[
 (K_{X_i}+t_i\alpha_i)\cdot C_i
       \geq-2d+t_i a_{ij}/l_{ij}>0,
\]
a contradiction. Therefore $\alpha_i\cdot R_i=0$.

Apply the ordinary rational step theorem to $(X_i,0)$ and $R_i$.
Its supporting nef class has the form $c_1(K_{X_i})+\kappa_i$
with $\kappa_i$ K\"ahler. Hence $\mathsf B_d$ supplies its
contraction, and the ordinary relative-positivity and flip arguments
give the projective contraction and its ordinary flip when needed.
All spaces remain compact K\"ahler, globally strongly
$\mathbb Q$-factorial, and klt. The preceding NQC transport applies.

It remains to check the positivity hypotheses at the next stage.
For a birational step, on a common resolution write
\[
 p^*K_{X_i}=q^*K_{X_{i+1}}+F_i,
 \qquad F_i\geq0,
 \qquad p^*\alpha_i=q^*\alpha_{i+1}.
\]
Here $F_i$ is the actual ordinary discrepancy divisor. Thus
\[
                 q^*Q_{i+1}=p^*Q_i+[F_i]
\]
is big, so $Q_{i+1}$ is big. Equality of the $\alpha$-pullbacks
also preserves non-bigness. We can therefore recompute the NQC
denominators and choose a fresh $t_{i+1}$ by
\eqref{repair:eq:adaptive-nqc-bound}. In the non-klt application,
$Q_i=B_i+\mathbf M_{X_i}$; the equality of the $\alpha$-pullbacks
preserves the original generalized adjoint and its b-datum.

If this procedure never reached a fibre-type contraction, it would
give an infinite sequence of ordinary rational klt $K$-negative
birational steps. Divisorial steps are finite by the strict decrease
of the span of analytic divisor classes, while small steps preserve
that span, as in the cycle-rank argument of Lemma~\ref{lem:cycle-loss}.
The remaining infinite flip tail would contradict ordinary rational
termination (Theorem~\ref{thm:gklt-three}). Hence it reaches an
ordinary Mori fibre space after finitely many steps. The final
$\alpha$-class descends to a nef weakly NQC class on its base by
the fibre-type argument above. No uniform bound for the $l_{ij}$,
and no application of a section theorem for rationally connected
fibres, is required.
\end{proof}


\subsubsection{Selected lower-dimensional good models}

\begin{lemma}[A selected good model in dimension at most three]
\label{repair:lem:lower-dimensional-good-model}
Fix $1\leq d\leq3$. Assume the generalized analytic cone theorem in
dimension $d$, the semiampleness assertion $B_d$, and termination of
ordinary rational klt programs of dimension at most three. Here $B_d$
means that a nef generalized klt adjoint with modified-big
boundary-plus-nef part is the pullback of a K\"ahler class under a
Moishezon contraction to a normal compact K\"ahler space. Its conclusion does not
include projectivity of that contraction. For ordinary rational
termination one may use Theorem~\ref{thm:gklt-three}, with zero nef
b-divisor.

Let $(X,B+\mathbf M)$ be a generalized klt pair on a normal compact
K\"ahler space of dimension $d$. Suppose that
\[
 A=K_X+B+\mathbf M_X
\]
is pseudoeffective and that $B+\mathbf M_X$ is modified big.
Then it has a good minimal model
$\phi:X\dashrightarrow X^{\rm m}$, where $X^{\rm m}$ is compact
K\"ahler and globally strongly $\mathbb Q$-factorial. The map extracts
no divisor. It can be selected so that every class in
$H^{1,1}_{\rm BC}(X)$ has a forward trace on $X^{\rm m}$.
No finiteness theorem for nearby generalized pairs is needed.
\end{lemma}

\begin{proof}
We first record precisely the elementary positivity input used in
passing to a smooth model. The modified-bigness resolution lemma
\cite[Lemma~2.5]{HLX2026} gives a projective smooth carrier
$\pi:Y\to X$ and an effective $\pi$-exceptional divisor $P$ such that,
on writing
\[
 K_Y+B_Y+M_Y=\pi^*A,
\]
the class $B_Y^++M_Y+\varepsilon P$ is big for every
$\varepsilon>0$. Here $M_Y$ is nef and $B_Y=B_Y^+-B_Y^-$, with
disjoint effective positive and negative parts. We can replace $P$ by
an effective divisor whose support contains every $\pi$-exceptional
prime. This positivity lemma follows directly from the definition of
modified bigness, the support theorem for a modification, and the
convexity of the big cone; it uses no contraction or termination theorem.
All the modifications here can be chosen projective.

For clarity, if $X$ is globally strongly $\mathbb Q$-factorial, that
input also has the following direct verification. The class
$C=B+\mathbf M_X=A-K_X$ is a big Bott--Chern class. Put
$Q=B_Y^++M_Y$, which is pseudoeffective. The difference
\[
 J=\pi^*C-Q=(K_Y-\pi^*K_X)-B_Y^-
\]
is an actual $\pi$-exceptional real divisor class. Given a positive
divisor $P$ supported on all exceptional primes, choose $s>0$ small
enough that $\varepsilon P-sJ\geq0$. Then
\[
 Q+\varepsilon P
  =(1-s)Q+s\pi^*C+(\varepsilon P-sJ)
\]
is big. This uses only that pullback preserves bigness and that adding
a pseudoeffective class to a big class preserves bigness.

We choose $\pi$ to resolve the supports just used. Choose
$\varepsilon>0$ small enough that
$\Delta=B_Y^++\varepsilon P$ has all coefficients below one.
Its support is simple normal crossing, so $(Y,\Delta)$ is klt and
\begin{equation}\label{repair:eq:good-model-first-error}
 K_Y+\Delta+M_Y=\pi^*A+F_Y,
 \qquad F_Y=B_Y^-+\varepsilon P\geq0.
\end{equation}
The support of $F_Y$ contains every $\pi$-exceptional prime, and
$\Delta+M_Y$ is big.

\smallskip\noindent
\emph{Smooth reduction with a K\"ahler nef part.}
Regularize a K\"ahler current in $\Delta+M_Y$ and resolve its analytic
singularities. This gives a projective modification $q:U\to Y$,
with $U$ smooth and compact K\"ahler, and
\[
 q^*(\Delta+M_Y)=\omega+[G],\qquad G\geq0,
\]
where $\omega$ is a K\"ahler class. To obtain a K\"ahler class rather
than a semipositive pullback, subtract a sufficiently small effective
exceptional class from the smooth part and add it to $G$. This is the
usual K\"ahler-current proof of the analytic Kodaira decomposition.
We resolve the supports involved at the same time. Write
$K_U=q^*K_Y+K_{U/Y}$; the divisor $K_{U/Y}$ is effective and
$q$-exceptional. For sufficiently small $\delta>0$, set
\[
 \beta=(1-\delta)q^*M_Y+\delta\omega,
 \qquad
 \Gamma_\delta=(1-\delta)q^*\Delta+\delta G-K_{U/Y}.
\]
The class $\beta$ is K\"ahler. At $\delta=0$ the second expression
is the crepant boundary of the klt pair $(Y,\Delta)$. Since its
support is fixed and finite, for small $\delta$ all coefficients of
$\Gamma_\delta$ remain below one. Its negative components are
$q$-exceptional. If $P_q$ is supported positively on every
$q$-exceptional prime, choose $\eta>0$ small enough that
\[
 \Gamma=\Gamma_\delta^++\eta P_q
\]
still has all coefficients below one. Its support is simple normal
crossing, so $(U,\Gamma)$ is klt. Put $\mu=\pi q$. We obtain
\begin{equation}\label{repair:eq:good-model-smooth-error}
 \widetilde A:=K_U+\Gamma+\beta
   =\mu^*A+E_\mu,
 \qquad
 E_\mu=q^*F_Y+\Gamma_\delta^-+\eta P_q\geq0.
\end{equation}
The support of $E_\mu$ contains every $\mu$-exceptional prime.
In particular $\widetilde A$ is pseudoeffective.

Choose a rational effective boundary $B_0$ sufficiently close to
$\Gamma$, on the same positive support, and put
\[
 H=\beta+[\Gamma-B_0],\qquad D=K_U+B_0.
\]
Openness of the K\"ahler cone makes $H$ K\"ahler, and $(U,B_0)$
remains klt. Fix a K\"ahler form representing $H$ and use its
pullbacks as one fixed nef b-datum. Thus
\begin{equation}\label{repair:eq:good-model-rationalization}
 \widetilde A=D+H,
\end{equation}
where $D$ is an ordinary rational klt adjoint. This auxiliary nef
b-datum need not equal the original $\mathbf M$.

\smallskip\noindent
\emph{Scaling only above one.}
Choose $T>1$ such that $D+TH$ is nef. Run the ordinary $D$-program
with scaling of this fixed $H$, stopping as soon as the forward trace
of $D+H$ is nef. Write its traces as $D_i,H_i$ and its successive
thresholds as $\lambda_i$, with $\lambda_{-1}=T$. The construction
of Proposition~\ref{prop:scaling-construction} applies in dimension
$d$: at a positive threshold it uses the generalized cone theorem,
finiteness of the negative rays for the big
boundary-plus-nef part, and the ordinary rational step proposition.
The latter uses precisely $B_d$ together with the integral descent
and relative-positivity arguments of Section~\ref{sec:absolute}.

More explicitly, if $D_i+H_i$ is not nef, convexity of the nef cone
gives $\lambda_i>1$, and there is an analytic extremal ray $R_i$
such that
\[
 D_i\cdot R_i<0,\qquad
 (D_i+\lambda_iH_i)\cdot R_i=0,\qquad H_i\cdot R_i>0.
\]
For the existence of a ray at the threshold, take parameters $t<\lambda_i$
increasing to $\lambda_i$. On a ray negative for $D_i+tH_i$,
nefness of the preceding upper-threshold class implies $H_i\cdot R>0$.
Thus, once $t>\lambda_i/2$, that ray is also negative for
$D_i+(\lambda_i/2)H_i$. The latter has only finitely many negative
rays. Its boundary-plus-nef trace is big, and
the pair is generalized klt by the interpolation argument below.
Every step performed is therefore negative for $D_i+H_i$.

Each birational step is an ordinary rational klt step contracting a
single analytic ray. It is projective on both sides, stays in the
compact K\"ahler globally strongly $\mathbb Q$-factorial category,
and has the forward class transport of
Section~\ref{repair:sec:real}. If $F_i\geq0$ is its ordinary
discrepancy divisor on a common resolution, crepancy at the scaling
threshold gives the discrepancy divisor
\begin{equation}\label{repair:eq:good-model-scaling-discrepancy}
                  (1-1/\lambda_i)F_i\geq0
\end{equation}
for $D_i+H_i$. More generally the divisor at parameter
$0\leq t\leq\lambda_i$ is $(1-t/\lambda_i)F_i$.
This proves the required generalized klt induction and makes the
next upper-threshold class nef. It also shows that
$q_i^*H_{i+1}=p_i^*H_i+F_i/\lambda_i$ on the common resolution;
in particular the forward traces $H_i$ stay big.

A Mori fibre step cannot occur while $\lambda_i>1$. Indeed the
pseudoeffective class $D_i+H_i$ would be negative on every curve in
a positive-dimensional fibre. Restrict a positive current in this
class to a general projective fibre, and then to a smooth resolution
of that fibre. Such a restriction is defined for a general fibre by
the local-potential definition and Fubini's theorem. Its intersection
with a power of an ample class is nonnegative, whereas a general
complete-intersection curve has class in $R_i$ and gives a negative
intersection. This is a contradiction. Pseudoeffectivity throughout
the program follows either by pushing forward positive currents, or
from the effective exceptional comparison and
Lemma~\ref{lem:positive-descent}.

The program cannot have infinitely many birational steps, by ordinary
rational termination in dimension at most three. If $D_i$ becomes
nef or a threshold reaches $[0,1]$, the convex interval between it
and the preceding nef upper-threshold class contains $D_i+H_i$.
Consequently the stopping rule produces a model $X^{\rm m}$ with
\[
                       A_{\rm m}=D_{\rm m}+H_{\rm m}\quad\text{nef}.
\]
Every step actually performed had $\lambda_i>1$.

\smallskip\noindent
\emph{Removing all the divisors introduced by the resolution.}
Take a common projective resolution with maps
$r:W\to U$, $v:W\to X^{\rm m}$ and $s=\mu r:W\to X$.
The comparisons \eqref{repair:eq:good-model-scaling-discrepancy} accumulate
to a Bott--Chern comparison with an actual effective divisor
\[
 r^*\widetilde A=v^*A_{\rm m}+[E],
 \qquad E\geq0,\quad E\text{ is }v\text{-exceptional}.
\]
Its coefficient is positive on the strict transform of every prime
on $U$ contracted by the program. Set
\[
                       G_0=r^*E_\mu-E.
\]
By \eqref{repair:eq:good-model-smooth-error},
$[G_0]=v^*A_{\rm m}-s^*A$. Hence $G_0$ is $s$-nef and
$s_*G_0=-s_*E\leq0$. The negativity lemma, applied to $-G_0$,
gives $G_0\leq0$. Thus
\begin{equation}\label{repair:eq:good-model-original-comparison}
 s^*A=v^*A_{\rm m}+[E-r^*E_\mu],
 \qquad E-r^*E_\mu\geq0.
\end{equation}
Since $E$ is $v$-exceptional and $E_\mu$ has positive coefficient on
every $\mu$-exceptional prime, every such prime was contracted.
The induced map $\phi:X\dashrightarrow X^{\rm m}$ therefore
extracts no divisor. The error in
\eqref{repair:eq:good-model-original-comparison} is $v$-exceptional and
has positive coefficient on each $\phi$-exceptional prime of $X$:
the subtracted term has coefficient zero on those primes.

We now return to the original nef b-datum $\mathbf M$. Keep it on
the same common carrier $W$. The generalized boundary defining the
original pair on $W$ is changed by subtracting the actual effective
divisor $E-r^*E_\mu$ in
\eqref{repair:eq:good-model-original-comparison}. It pushes forward to
$\phi_*B$ because $\phi$ extracts no divisor. It defines
$K_{X^{\rm m}}+\phi_*B+\mathbf M_{X^{\rm m}}=A_{\rm m}$, and
its coefficients are below one because the original pair is gklt.
Thus $X^{\rm m}$ is a minimal model of the original generalized
pair, with its original b-datum. No assertion that the original
boundary remains effective or generalized klt on the intermediate
auxiliary models is required.

Finally the auxiliary generalized pair
$(X^{\rm m},B_{0,\rm m}+\overline H)$ is gklt, and its
boundary-plus-nef trace is big. Apply $B_d$ to this pair to see that
$A_{\rm m}$ is semiample, so the model is good. For any original
class $\theta\in H^{1,1}_{\rm BC}(X)$, start with $\mu^*\theta$
and use the forward transport across each ordinary analytic-ray
step. This defines a class $\theta_{\rm m}$ on $X^{\rm m}$ and
an actual exceptional-divisor comparison on $W$. In particular
\eqref{repair:eq:good-model-original-comparison} identifies the forward
trace of $A$ with $A_{\rm m}$. Only this forward transport is
asserted; no surjectivity from all Bott--Chern classes on the final
model is used.
\end{proof}


\subsection{Negative parts and lower-dimensional good models}
\label{repair:sec:negative-part-replacement}

All equalities between $(1,1)$-classes in this section are equalities in
Bott--Chern cohomology. Equalities between divisors are stated separately.
The spaces are normal compact K\"ahler spaces with rational singularities;
resolutions and the morphisms called projective carry actual relatively
ample line bundles. We write $N_\sigma(\xi)$ for the divisorial negative
part of a pseudoeffective class. On a singular space it is defined by
pushing forward the negative part on a resolution.

We use the following properties of the negative part:
\begin{enumerate}
\item A nef class has zero negative part. If $\xi=P+[F]$, with $P$ nef
and $F\geq0$ an effective real Cartier divisor, then
$N_\sigma(\xi)\leq F$.
\item If $r:W\to T$ is a modification, $\xi$ is pseudoeffective, and
$E\geq0$ is $r$-exceptional, then
\[
 N_\sigma(r^*\xi+[E])=N_\sigma(r^*\xi)+E,
 \qquad r_*N_\sigma(r^*\xi)=N_\sigma(\xi).
\]
\item Minimal multiplicities are subadditive and homogeneous. In
particular, adding a nef class can only decrease them. For a K\"ahler
class $\omega$,
\[
 \nu(\xi,Q)=\lim_{\varepsilon\downarrow0}
                 \nu(\xi+\varepsilon\omega,Q).
\]
On a big class the minimal multiplicity is the infimum of the generic
Lelong numbers of its closed positive currents.
\end{enumerate}
These are the usual properties of the divisorial Zariski decomposition
\cite[\S3]{Boucksom2004}. For the exceptional-divisor identity and its
extension to normal spaces, only
\cite[Appendix~A, Lemmas~A.3, A.5 and A.7]{DHY2023} are needed.
No assertion about transporting arbitrary Bott--Chern classes across a
small bimeromorphic map is used here.

\subsubsection{The initial relative program}

\begin{lemma}[The absolute negative part after a relative program]
\label{repair:lem:negative-part-relative-program}
Let $\mu:U\to X$ be a projective resolution, let $\alpha$ be nef on $X$,
let $c>0$, and let $F_U\geq0$ be $\mu$-exceptional. Put
\[
 D_U=c\mu^*\alpha+[F_U].
\]
Suppose that $\phi:U\dasharrow V$ is a finite $D_U$-negative program,
possibly relative to another base, and denote its transformed class and
divisor by $D_V$ and $F_V=\phi_*F_U$. Then
\[
                         N_\sigma(D_V)=F_V.
\]
\end{lemma}

\begin{proof}
Take a common projective resolution $p:W\to U$, $q:W\to V$.
The negativity comparison for the finite program is
\[
                p^*D_U=q^*D_V+[E],
 \qquad E\geq0,\quad E\text{ is }q\text{-exceptional}.
\]
This comparison is valid for a relative negative program: its proof
uses the negativity of the contracted rays and the positive adjoint on
each flipped side, not absolute nefness of the final adjoint.
Since $p^*F_U$ is effective and exceptional over $X$, the exceptional
identity and nefness of $(\mu p)^*\alpha$ give
\[
 N_\sigma(p^*D_U)=p^*F_U
                 =N_\sigma(q^*D_V)+E.
\]
Pushing forward by $q$ proves the assertion. In particular, the
negative part in this statement is absolute, even when the program
that produced $V$ was relative.
\end{proof}

\subsubsection{Descent of the effective vertical divisor}

\begin{lemma}[Vertical numerical triviality]
\label{repair:lem:negative-part-vertical-descent}
Let $g:V\to S$ be a projective surjective morphism with connected fibres.
Assume that $S$ is globally Weil $\mathbb{Q}$-factorial. Let $F\geq0$ be
an effective real Cartier divisor on $V$, vertical over $S$, such that
\[
                    F\cdot C=0
       \quad\text{for every curve }C\text{ contracted by }g.
\]
Then there is an effective real Cartier divisor $F_S$ on $S$ such that,
as actual divisors,
\[
                              F=g^*F_S.
\]
\end{lemma}

\begin{proof}
For a prime divisor $P\subset S$, write $m_Q$ for the multiplicity of
$Q$ in $g^*P$, where $Q$ ranges over the prime divisors dominating $P$.
These multiplicities may be computed over the smooth generic locus of
$P$, where $P$ is Cartier. Put
\[
 b_P=\min_{g(Q)=P}\frac{\operatorname{coeff}_Q F}{m_Q},
 \qquad F_S=\sum_P b_PP.
\]
Only finitely many $b_P$ are nonzero, since any such $P$ is the image of
a component of $F$. Thus $F_S$ is a genuine effective Weil real divisor.
Global Weil $\mathbb{Q}$-factoriality makes this finite divisor real
Cartier.

We first check equality over codimension one in $S$. Work near a
general point of $P$ and restrict to a transverse disk. Resolving the
source, and cutting by general relative ample hypersurfaces if the
relative dimension exceeds one, reduces to a projective surface over
that disk with connected fibres. These operations can be performed
over a relatively compact Stein neighbourhood; they do not require
global sections on $S$. The intersection matrix of the components of
the special fibre is negative semidefinite, with kernel generated by
the full fibre with its multiplicities. The restriction of $F$ has
zero intersection with every fibre component. Its coefficients are
therefore proportional to those multiplicities. Equivalently, all the
ratios in the definition of $b_P$ are equal. When $g$ is birational this
assertion over the generic point of $P$ is immediate.

Consequently
\[
                              G=F-g^*F_S
\]
is a signed real Cartier divisor supported over a subset of codimension
at least two in $S$. Moreover, $G$ is numerically trivial over $S$.
We check that such a signed exceptional divisor is zero. The assertion
is local on $S$. Over a relatively compact Stein neighbourhood, take
$d=\dim V-\dim S$ general relative ample hypersurfaces and resolve their
intersection. They give a projective generically finite morphism
$H\to S$. The cuts may be chosen to meet any prescribed component of
$G$ in a nonzero divisorial trace. After normalization and Stein
factorization, $H\to S$ factors as a projective birational morphism
$H\to S'$ followed by a finite morphism $S'\to S$. The restricted
divisor $G|_H$ is exceptional for $H\to S'$ and numerically trivial
over $S'$. Applying the ordinary exceptional negativity lemma to both
$G|_H$ and $-G|_H$ gives $G|_H=0$. The choice of the cuts therefore
excludes every nonzero component of $G$. Hence $G=0$.

The surface argument above is also the usual proof of the
degenerate-divisor negativity statement: after subtracting the minimum
multiple of the full fibre, an effective residual fibre divisor omits
a component and cannot be numerically trivial. Notice that projectivity
of $g$ was a hypothesis; it was not deduced from factoriality.
\end{proof}

\subsubsection{The negative part on the lower-dimensional base}

\begin{lemma}[Detecting the negative part by pullback currents]
\label{repair:lem:negative-part-base}
Let $g:V\to S$ be a surjective morphism. Suppose that $\alpha_S$ is nef,
$F_S\geq0$ is real Cartier, $c>0$, and
\[
 D_S=c\alpha_S+[F_S],\qquad D_V=g^*D_S,\qquad
 N_\sigma(D_V)=g^*F_S.
\]
Then $N_\sigma(D_S)=F_S$.
\end{lemma}

\begin{proof}
Nefness gives $N_\sigma(D_S)\leq F_S$. Fix a prime divisor $P\subset S$
and a prime divisor $Q\subset V$ dominating it. Write
$m=\operatorname{mult}_Q(g^*P)>0$ and $b=\operatorname{coeff}_P F_S$.
At their generic smooth points, pullback of positive currents satisfies
\[
                         \nu(g^*T,Q)=m\nu(T,P).
\]
Indeed, the divisorial term $\nu(T,P)[P]$ pulls back with multiplicity
$m$, and the residual current has zero generic Lelong number there.

Choose K\"ahler forms $\omega_S,\omega_V$ and a constant $C>0$ for which
$C\omega_V-g^*\omega_S$ is K\"ahler. For $\varepsilon>0$, let $T_\varepsilon$
be any positive current in the big class $D_S+\varepsilon[\omega_S]$.
Minimality of the multiplicity and monotonicity under adding a nef
class give
\begin{align*}
 m\nu(T_\varepsilon,P)
 &=\nu(g^*T_\varepsilon,Q)\\
 &\geq\nu(D_V+\varepsilon g^*[\omega_S],Q)\\
 &\geq\nu(D_V+C\varepsilon[\omega_V],Q).
\end{align*}
Taking the infimum over $T_\varepsilon$, and then letting
$\varepsilon\downarrow0$, gives
\[
 m\nu(D_S,P)\geq\nu(D_V,Q)
                 =\operatorname{coeff}_Q(g^*F_S)=mb.
\]
This proves the reverse inequality for every $P$.
All multiplicities are unchanged by resolving away from the generic
points in question, so the same argument applies to the normal spaces
under consideration.

The $\varepsilon$ perturbation is necessary: at a pseudoeffective
boundary class, the infimum over its exact positive currents need not
equal its minimal multiplicity. The argument uses that equality only
for the big perturbed classes.
\end{proof}

\begin{lemma}[Pullback when the positive part is nef]
\label{repair:lem:negative-part-nef-pullback}
Suppose
\[
 \xi=P+[F],\qquad P\text{ nef},\quad F\geq0\text{ real Cartier},
 \qquad N_\sigma(\xi)=F.
\]
For every projective resolution $p:W\to S$,
\[
                         N_\sigma(p^*\xi)=p^*F.
\]
\end{lemma}

\begin{proof}
Write $N'=N_\sigma(p^*\xi)$. Since $p^*P$ is nef, $N'\leq p^*F$.
Birational invariance gives $p_*N'=F$, so
$E=p^*F-N'\geq0$ is $p$-exceptional. The positive part of $p^*\xi$ is
\[
                         p^*P+[E],
\]
and is modified nef. If $E\neq0$, exceptional negativity in its
covering-curve form provides a component of $E$ covered by curves $C_t$
contracted by $p$, with $E\cdot C_t<0$
\cite[Appendix~A, Lemma~A.3]{DHY2023}. A modified nef class has nonnegative
intersection with a general member of a family of curves covering a
prime divisor: use regularization to choose currents with analytic singularities,
arbitrarily small negative part, and zero generic divisorial Lelong number;
then restrict to a general member and let the negative bound tend to zero. But here
\[
                       (p^*P+E)\cdot C_t=E\cdot C_t<0,
\]
a contradiction. Thus $E=0$.
\end{proof}

\begin{proposition}[Any lower-dimensional good model suffices]
\label{repair:prop:negative-part-good-model}
Suppose the initial relative program has the data in
Lemma~\ref{repair:lem:negative-part-relative-program}. Assume in addition that $F_V$ is a real Cartier divisor, that
$\alpha_V$ is nef and agrees with $\mu^*\alpha$ on a common
resolution, and that there is an already projective connected-fibre
morphism $g:V\to S$ with
\[
\alpha_V=g^*\alpha_S,\qquad D_V=g^*D_S.
\]
Suppose there is a projective small modification
$s:S^{\mathrm q}\to S$ such that $S^{\mathrm q}$ is globally Weil
$\mathbb Q$-factorial; the identity is allowed. Assume that the
lower-dimensional adjoint $s^*D_S$ has a good model:
there are a normal compact K\"ahler space $S_m$, a common projective
resolution $p:W\to S^{\mathrm q}$, $q:W\to S_m$, and a nef semiample
class $D_m$ such that
\[
                  p^*s^*D_S=q^*D_m+[E],
 \qquad E\geq0,\quad E\text{ is }q\text{-exceptional}.
\]
Then $\alpha$ is semiample on $X$. If the contraction defining
semiampleness of $D_m$ is Moishezon, the resulting contraction of
$\alpha$ is Moishezon as well. In particular, it is unnecessary to
lift a program on $S$ to $V$, or to require that a chosen program
producing $S_m$ preserve $\alpha_S$ at every intermediate step.
\end{proposition}

\begin{proof}
Nefness descends under the surjective morphism $g$, so $\alpha_S$ is nef.
The class of $F_V$ is pulled back from $S$, so $F_V$ is numerically
trivial over $S$. It is therefore vertical: an effective horizontal
component would restrict to a nonzero effective divisor on a general
projective fibre, with positive intersection against a suitable power
of an ample class, contradicting numerical triviality. For a
birational $g$, verticality is automatic by dimension.
First take a projective resolution $\pi:\widetilde V\to V$ of the main
component of $V\times_S S^{\mathrm q}$. The induced map
$\widetilde g:\widetilde V\to S^{\mathrm q}$ is projective and has
connected fibres. Pull back $\alpha_V,D_V,F_V$ along $\pi$, and pull
back $\alpha_S,D_S$ along $s$. All displayed class identities are
preserved, and $\pi^*F_V$ remains an effective vertical divisor.
Lemma~\ref{repair:lem:negative-part-relative-program} gives
$N_\sigma(D_V)=F_V$ before this modification. Because
$D_V=c\alpha_V+[F_V]$ with $\alpha_V$ nef,
Lemma~\ref{repair:lem:negative-part-nef-pullback} then gives
\[
              N_\sigma(\pi^*D_V)=\pi^*F_V.
\]
We may therefore replace $(V,S,g)$ by
$(\widetilde V,S^{\mathrm q},\widetilde g)$ and suppress the new
superscripts in the rest of the proof. The good-model comparison now
reads $p^*D_S=q^*D_m+[E]$. No assertion that the crepant subboundary on
$\widetilde V$ is effective is needed: this space is used only for
classes, currents and the effective divisor $\pi^*F_V$. The
lower-dimensional generalized pair is the crepant pullback to the
small model $S^{\mathrm q}$.

The class identity $[F_V]=g^*(D_S-c\alpha_S)$ shows that $F_V$ is
numerically trivial over $S$. Lemma~\ref{repair:lem:negative-part-vertical-descent}
gives an actual effective real Cartier divisor $F_S$ with
$F_V=g^*F_S$. Injectivity of pullback yields
\[
                         D_S=c\alpha_S+[F_S].
\]
The established identity $N_\sigma(D_V)=F_V$ and
Lemma~\ref{repair:lem:negative-part-base} therefore give
$N_\sigma(D_S)=F_S$, and
Lemma~\ref{repair:lem:negative-part-nef-pullback} gives
\[
                         N_\sigma(p^*D_S)=p^*F_S.
\]
On the other hand $q^*D_m$ is nef, so the exceptional-divisor identity
applied to the good-model comparison gives
\[
                         N_\sigma(p^*D_S)=E.
\]
Thus $E=p^*F_S$ as actual divisors. Subtracting them from the comparison
proves the exact class equality
\[
                         q^*D_m=c\,p^*\alpha_S.
\]
This proves the needed preservation of $\alpha_S$ from the final good
model, without an assumption about its intermediate models.

Choose a contraction $h:S_m\to T$ to a normal compact K\"ahler space
and a K\"ahler class $\kappa$ on $T$ with $D_m=h^*\kappa$. Taking a
resolution of the main component of $V\times_S W$, and then a common
projective resolution with $U$, produces a projective modification
$r:R\to X$ and a holomorphic map $\ell:R\to T$ satisfying
\[
                         r^*\alpha=\frac1c\ell^*\kappa.
\]
For every curve $C$ in a fibre of $r$ this equality implies
$\ell(C)$ is a point, since a K\"ahler class has positive degree on
every nonconstant image curve. The fibres of the projective modification
$r$ are connected projective complex spaces, hence are connected by
chains of curves. Therefore $\ell$ is
constant on each fibre of $r$. The factorization theorem for a proper
map onto a normal space gives a holomorphic map $f:X\to T$ with
$\ell=f\circ r$. Pushing forward the class equality by $r$ gives
\[
                              \alpha=f^*(\kappa/c).
\]
The maps from the main fibre-product component to $S_m$ have connected
general fibres; their Stein factorizations are finite birational over
the normal space $S_m$, hence have connected fibres everywhere.
Together with the connected fibres of $h$, this shows that $\ell$,
and therefore $f$, has connected fibres. This is the required
semiampleness of $\alpha$.

Finally, $R\to S_m$ is projective: it is obtained from the projective
map $g$ by base change, followed by projective resolutions and the
projective map $q$. If $h$ is Moishezon, choose a projective surjection
$H\to S_m$ such that $H\to T$ is projective. A component of
$R\times_{S_m}H$ dominating $R$ is projective and surjective over $X$,
and its map to $T$ is projective. This is a Moishezon witness for $f$.
\end{proof}


\subsection{Closing the contraction induction}
\begin{proposition}[The restricted induction]
\label{repair:prop:restricted-bpf-induction}
Let $\mathsf T_d$ denote termination of ordinary rational klt
programs in dimension at most $d$. The following implications hold:
\begin{align*}
 \mathsf C_{d-1}&\Longrightarrow\mathsf C_{d,\mathrm{big}},\\
 \mathsf C_{d,\mathrm{big}}+\mathsf B_{d-1}+\mathsf G_{d-1}
   &\Longrightarrow\mathsf B_d \qquad (2\leq d\leq4),\\
 \mathsf B_d+\mathsf T_d
   &\Longrightarrow\mathsf G_d \qquad (d\leq3),\\
 \mathsf B_d+\mathsf C_{d-1}+\mathsf T_d
   &\Longrightarrow\mathsf C_d \qquad (d\leq3).
\end{align*}
Consequently $\mathsf B_4$ holds. No assertion $\mathsf G_4$ or
$\mathsf C_4$ is required.
\end{proposition}

\begin{proof}
We specify the inputs from \cite{HX2026,HLX2026} and the arguments
used below. The inputs are the cone theorem, projective
relative vanishing and minimal-model theory, projective small
factorializations and dlt modifications, the projective canonical
bundle formula, and the nef-and-big K\"ahler criterion. None is
used to promote an arbitrary map from a globally strongly
$\Q$-factorial space to a projective map.

\smallskip\noindent
\emph{The big non-klt contraction.}
Apply the construction of \cite[Section~3]{HX2026}. On its smooth
modification $\nu:U\to X$ it writes
\[
 \nu^*\alpha=[D_U]+\eta_U,
 \qquad D_U\geq0,\quad \eta_U\text{ K\"ahler}.
\]
The induction over the jumping coefficients of the multiplier
ideal reduces the new reduced stratum to dimension at most $d-1$.
The given contraction on the old non-klt subspace and
$\mathsf C_{d-1}$ provide the Moishezon maps of those strata.
Every map that is promoted to a projective map in that construction
contracts exactly the curves on which $\nu^*\alpha$ vanishes.
On each such curve the actual real line bundle $-D_U$ has degree
$\eta_U\cdot C$. On a further projective common resolution
$q:U''\to U$, choose an effective exceptional divisor $E$ with
$-E$ relatively ample and choose $\varepsilon>0$ sufficiently small.
Use the decomposition
\[
 q^*\nu^*\alpha=[D'']+\eta'',\qquad
 D''=q^*D_U+\varepsilon E,\qquad
 \eta''=q^*\eta_U-\varepsilon[E].
\]
Here $\eta''$ is K\"ahler, and the actual real line bundle $-D''$
has degree $\eta''\cdot C$ on every contracted curve. Restrict this
decomposition to the smooth reduced components in the gluing
construction. Lemma~\ref{repair:lem:actual-line-polarization} then supplies
precisely the relative ampleness needed in place of
\cite[Lemma~2.42]{HX2026}, including the common-resolution step in its
Claim~3.5. The uncorrected pullback $q^*\eta_U$ is not being treated
as a K\"ahler class.

The exact sequences of multiplier ideals, relative vanishing,
finite pushouts, and extension from sufficiently thickened
$\operatorname{Supp}D_U$ in that proof then apply unchanged.
They give the birational Moishezon contraction in
$\mathsf C_{d,\mathrm{big}}$. In particular the gluing step uses
neither generalized termination nor a projectivity criterion for
an undetected ray.

\smallskip\noindent
\emph{Preparation for both cases of $\mathsf B_d$.}
The modified-big perturbation
\cite[Lemma~2.23]{HX2026} supplies a nef datum that dominates a
K\"ahler class on its carrier. Subtracting a sufficiently small
multiple of the pullback of a K\"ahler class $\omega_X$ still leaves
a nef datum on that carrier. Thus
$\alpha$ is a gklt adjoint plus $\varepsilon\omega_X$, and the cone
argument of \cite[Lemma~2.45]{HX2026} makes $\alpha$ NQC.
This preparation applies whether or not $\alpha$ is big.

\smallskip\noindent
\emph{The nef-and-big case of $\mathsf B_d$.}
For a gklt pair the multiplier ideal is the unit ideal, so the
preceding big non-klt assertion now applies and gives a birational
contraction $h:X\to Y$.
On a projective resolution $p:W\to X$ chosen projective over $Y$,
relative Kawamata--Viehweg vanishing for the effective exceptional
divisor $-\lfloor B_W\rfloor$ gives
$R^ih_*\mathcal O_X=0$ for $i>0$. Thus $Y$ has rational
singularities. Proper birational Bott--Chern descent
\cite[Lemma~8.7]{DHP2024} gives $\alpha=h^*\gamma$.
The descended adjoint is gklt, nef and big, and has no trivial
curve. The criterion \cite[Theorem~4.3]{HLX2026} makes $\gamma$
K\"ahler.

For completeness, the projectivity input in that criterion is
restricted to a map $S\to Z'$ between smooth compact K\"ahler
manifolds with rationally connected general fibre: $S$ is the
chosen smooth divisor on a resolution and $Z'$ is the resolution
of the null-locus component. The smooth-source, smooth-base
argument in \cite[Theorem~3.1, Step~1]{ClaudonHoring2024} applies. Its subsequent
relative program starts with this projective map. Thus no general singular-source projectivity assertion is needed
for this application of the criterion.

\smallskip\noindent
\emph{The non-big case of $\mathsf B_d$.}
Use a projective small factorialization and the modified-big
perturbation of the pair. The non-pseudo-effectivity of $K_X$
gives an MRC fibration by \cite{Ou2025}. Choose a smooth K\"ahler
MRC base $Z$ and a smooth modification $\mu:U\to X$ such that
$U\to Z$ is projective. Only the smooth case of
\cite[Theorem~3.1, Step~1]{ClaudonHoring2024} is needed: pullback identifies
the holomorphic two-forms since the general fibre is rationally
connected. The same rational Hodge correction identifies any
$(1,1)$-class on $U$, modulo a class pulled back from $Z$, with
an actual real line bundle.

Write, as in \cite[Theorem~4.1, Claim~4.1]{HX2026},
\[
 D_U(a)=K_U+\Delta_U+\mathbf M_U+a\alpha_U
     =(a+1)\alpha_U+F_U,
 \quad \alpha_U=\mu^*\alpha,\quad
 \Delta_U,F_U\geq0,
\]
where $\Delta_U$ is klt and $F_U$ is $\mu$-exceptional.
The modified-big replacement is chosen, as in the cited Claim~4.1,
so that the nef datum $\mathbf M_U$ on this smooth carrier is
K\"ahler. If a further projective resolution is needed, subtract a
sufficiently small effective exceptional divisor from its nef part
and add that divisor to the subboundary; this preserves the
adjoint and the klt inequalities.
Choose $a_0$ using the uniform bound of
Proposition~\ref{repair:prop:relative-nqc-uniform}. It makes the relative
$D_U(a_0)$-program $\alpha_U$-trivial and preserves that choice
through every step. This is a
\emph{projective} relative program from its first step: the nef
data are represented by real line bundles modulo $Z$.
Lemma~\ref{repair:lem:fixed-smooth-base-algebraicity} applies on every
intermediate model over this fixed smooth base.
More precisely, put $r:U\to Z$ and choose one actual global real
line bundle $L$ with
\[
 \mathbf M_U+a_0\alpha_U=c_1(L)+r^*\gamma.
\]
The bundle $L$ is relatively ample, because $\mathbf M_U$ is
K\"ahler and $\alpha_U$ is nef. On each of a fixed finite collection
of relatively compact Stein charts of $Z$, choose an effective
real divisor $\Theta$ representing $L$, with $(U,\Delta_U+\Theta)$
klt. The representatives can retain a positive relatively ample
part in their boundaries. Here the equality
$K_U+\Delta_U+\Theta\sim_{\R}K_U+\Delta_U+L$ is an actual
real linear equivalence of line bundles; the equality with
$D_U(a_0)$ modulo $r^*\gamma$ is separately an equality of
Bott--Chern classes.

For the ample-model comparison below, make the following additional
choices before running the program. Take the fixed finite cover in
the form $W_\ell\Subset Y_\ell\subset Z$, where $Y_\ell$ is a Stein
coordinate chart and $W_\ell$ is a closed coordinate polydisc, with
the interiors of the $W_\ell$ covering $Z$. Each $W_\ell$ is a
semianalytic Stein compact, so it satisfies condition~(P4) of
\cite{Fujino2022}; thus condition~(P) holds for the restriction of every
normal projective model over $Z$ to $Y_\ell$.

Write
\[
 \alpha_i=c_1(N_i)+f_i^*\Gamma
\]
for the real line bundles $N_i$ transported in
Proposition~\ref{repair:prop:relative-nqc-uniform}; put $N_i=0$ if
$\alpha_U=0$. These bundles are nef over $Z$ and pull back from
each contraction target, on both sides of a flip. Set $a_j=a_0+j$
for $j=1,2$. The real line bundles $L+jN_U$ are relatively ample.
On each initial chart choose, once and for all,
\[
 \Theta_j\geq0,\qquad \Theta_j\sim_{\R}L+jN_U,
 \qquad (U,\Xi_{j,U}:=\Delta_U+\Theta_j)\ \text{klt}.
\]
Such representatives are obtained from general sufficiently
divisible relative sections; their boundaries are relatively big
because $L+jN_U$ is relatively ample and $\Delta_U\geq0$.
With $\Xi_{0,U}:=\Delta_U+\Theta$, these are actual real linear
equivalences
\[
 K_U+\Xi_{j,U}\sim_{\R}K_U+\Xi_{0,U}+jN_U
\]
on each chart, separately from the Bott--Chern comparison with the
fixed base classes.

Let $\Xi_{j,i}$ denote their birational transforms. Pushforward of
these linear equivalences and the actual descent of $N_i$ give
\[
 K_{U_i}+\Xi_{j,i}
 \sim_{\R}K_{U_i}+\Xi_{0,i}+jN_i
\]
at every step. Thus the contracted adjoints have the same degrees
on the negative side, and the flipped adjoints have the same degrees
on the positive side. Every step of the chosen program is therefore
also a $(K_{U_i}+\Xi_{j,i})$-negative step. Ordinary discrepancy
negativity preserves klt for these fixed pairs. Their boundaries
remain effective, and relative bigness is preserved by these
birational pushforwards, as in the proof of
\cite[Lemma~11.16]{Fujino2022}.

Fix these ordinary pairs on the initial charts. Their actual
linear equivalences push forward throughout the program. On a
common resolution at each step, exceptional negativity preserves
the Bott--Chern comparison with $D_U(a_0)$ modulo the fixed base
class. Thus every subsequent contraction or flip is a step for
these fixed ordinary pairs. The projective ordinary theorems
\cite{Fujino2022,FujinoCompact2026} apply. In particular, the finiteness
of weak log canonical models in \cite[Theorem~E]{Fujino2022} does not
require local $\Q$-factoriality of each intermediate model. Applied
on the finite chart cover, its chamber argument gives termination
of the global projective program with scaling. Each small step
is constructed by its global rational detector algebra as above;
uniqueness of the relatively ample model identifies the local
ordinary flips on overlaps. This supplies a good
minimal model $U\dasharrow V$ over $Z$. Here one uses the projective case of
\cite[Proposition~2.30]{HX2026}, not its general part~(3).

Here is an explicit ample-model comparison that gives the required
descent of $\alpha_V$. Both $D_V(a_0)$ and $\alpha_V$ are nef over
$Z$. For $j=1,2$, the ordinary adjoint $K_V+\Xi_{j,V}$ is
$\R$-Cartier and represents
\[
 D_V(a_j)-f_V^*(\gamma+j\Gamma)
   =D_V(a_0)+j\alpha_V-f_V^*(\gamma+j\Gamma).
\]
It is consequently nef over every $W_\ell$. The pair
$(V,\Xi_{j,V})$ is klt and its effective boundary is relatively
big. Apply \cite[Lemma~11.15]{Fujino2022} with
$A=\Xi_{j,V}$ and $B=0$. Its condition on non-klt centres is
vacuous. After shrinking around $W_\ell$, it gives
\[
 \Xi_{j,V}\sim_{\R}A_j+B_j,\qquad
 A_j\geq0\ \text{relatively ample},\quad B_j\geq0,
 \qquad (V,A_j+B_j)\ \text{klt}.
\]
In particular this pair is log canonical, supplies the klt boundary
required in \cite[Theorem~8.3]{Fujino2022}, and its adjoint is
$\R$-Cartier and nef over $W_\ell$. That theorem makes
$K_V+\Xi_{j,V}$ semiample on a neighbourhood of $W_\ell$.
The pushed-forward actual linear equivalences identify these local
adjoints with restrictions of the same global real line-bundle
data. Their ample models therefore agree on overlaps by uniqueness
and glue as above. Hence $D_V(a_1)$ and $D_V(a_2)$ have projective
relative ample-model contractions over $Z$. Their zero curves are exactly the common zero
curves of $D_V(a_0)$ and $\alpha_V$. Their projective ample-model
contractions therefore have the same fibres: these fibres are
connected by curves, and each contraction is constant on the
fibres of the other. Identify the two normal targets, writing the
common contraction as $g:V\to S$. If
$D_V(a_j)=g^*\lambda_j$, with $\lambda_j$ relatively K\"ahler over
$Z$, subtraction gives
\[
 \alpha_V=g^*(\lambda_2-\lambda_1)=g^*\alpha_S.
\]
This is an equality of Bott--Chern classes, not merely of curve
degrees. Since the initial program is $\alpha_U$-trivial, it is
also $D_U(a_1)$-negative. From now on put $a=a_1$ and suppress
the argument in $D_U(a)$ and $D_V(a)$.

The maps $V\to Z$ and $S\to Z$ are projective. Moreover $g$ is
projective: a relatively ample bundle for $V\to Z$ restricts to
an ample bundle on every fibre of $g$, and hence is $g$-ample.
The argument of \cite[Theorem~4.1, Claim~4.2]{HX2026}, using
projective relative semipositivity, gives $\dim S<d$.
The projective canonical bundle formula
\cite[Theorem~2.53]{HX2026} provides
\[
 D_S=K_S+B_S+\mathbf N_S+a\alpha_S,
 \qquad D_V=g^*D_S,
\]
with a gklt pair on $S$ and modified-big nef data.
The class $\alpha_S$ is nef by descent of nefness through the
surjective projective map $g$.
Thus adding $a\overline{\alpha_S}$ to the nef datum preserves the
gklt condition and modified bigness: on a common carrier this
addition is a nef pullback and preserves the existing big class.

Take a projective small factorialization $s:S^{\mathrm q}\to S$
and resolve the main component of $V\times_S S^{\mathrm q}$.
The resulting maps $\pi:\widetilde V\to V$ and
$\widetilde g:\widetilde V\to S^{\mathrm q}$ are projective,
and $\widetilde V$ is compact K\"ahler. Pull back $D_V$,
$\alpha_V$ and $F_V$ to $\widetilde V$, and $D_S$ and
$\alpha_S$ to $S^{\mathrm q}$. The pair on $S^{\mathrm q}$ is
gklt and its boundary-plus-nef class remains modified big.
Lemma~\ref{repair:lem:negative-part-relative-program} first gives
$N_\sigma(D_V)=F_V$. Lemma~\ref{repair:lem:negative-part-nef-pullback} then gives
$N_\sigma(\pi^*D_V)=\pi^*F_V$. The new total space is used only as a
carrier for this equality and for pullbacks of positive currents;
no effective boundary on it is needed. Rename these spaces $V$
and $S$. No extra relative program from Claim~4.4 of the cited proof is needed.

The vertical-divisor descent and negative-part lemmas above now
give actual effective divisors $F_S,F_V$ with
\[
 F_V=g^*F_S,\qquad
 D_S=(a+1)\alpha_S+[F_S],\qquad N_\sigma(D_S)=F_S.
\]
Apply $\mathsf G_{d-1}$ to the generalized pair with adjoint $D_S$.
On a common resolution $p:W\to S$, $q:W\to S_m$ of the resulting
good model, write
\[
 p^*D_S=q^*D_m+[E],\qquad E\geq0
 \text{ exceptional over }S_m.
\]
The negative-part comparison proved above gives
$E=p^*F_S$, and hence
\[
 q^*D_m=(a+1)p^*\alpha_S.
\]
Since $D_m$ is semiample, $\alpha_S$ is semiample after descent
through $p$. Pulling back by $g$ and using the
$\alpha$-trivial initial program gives semiampleness of
$\alpha$ on $X$. The target is compact K\"ahler, and the
resulting map is Moishezon, as follows by taking common projective
modifications of the maps just constructed. This proves the required descent without lifting a lower-dimensional
program to $V$.

\smallskip\noindent
\emph{Selected good models in dimensions at most three.}
Lemma~\ref{repair:lem:lower-dimensional-good-model} gives
$\mathsf G_d$ from $\mathsf B_d$ and ordinary rational
termination. Its smooth input has the form
\[
 \mu^*A+[E_\mu]=K_U+B_q+H,
 \qquad B_q\text{ rational klt},\quad H\text{ K\"ahler},
\]
where $E_\mu$ is effective with all $\mu$-exceptional divisors in
its support. The ordinary $(K_U+B_q)$-program with scaling of $H$
is stopped at the first threshold at most one. Every performed
step is negative for the displayed adjoint. Termination is the
ordinary specialization of Theorem~\ref{thm:gklt-three},
with zero nef b-divisor. Its proof uses the projective relative and
discrepancy arguments of Section~\ref{sec:threefold}, not
$\mathsf B_d$ or $\mathsf G_d$.
The final nef adjoint is semiample by $\mathsf B_d$.
Negativity removes the added exceptional divisors and gives the
claimed model of $X$. All forward transforms used here are those
of the detected ordinary steps; no inverse identification of
Bott--Chern groups is needed.

\smallskip\noindent
\emph{The non-big non-klt contraction in dimensions at most three.}
Take the projective dlt modification used in
\cite[Section~6]{HX2026}; its source is globally strongly
$\Q$-factorial and its underlying space is klt. On this source
\[
 Q=\alpha-c_1(K_X)=B+\mathbf M_X
\]
is a Bott--Chern class. Its modified bigness therefore implies
bigness, as in \cite[Definition~2.8]{HX2026}. Apply
Lemma~\ref{repair:lem:adaptive-nqc-program}. It gives a finite ordinary
$K$-negative program, crepant for $\alpha$, ending in an
$\alpha$-trivial Mori fibre space $f:X'\to Z$. Each step exists
by $\mathsf B_d$ and is projective by its canonical rational
line-bundle detector. The auxiliary number $t_i$ is chosen anew
at each model; no uniform denominator claim is used.

All remaining steps of
\cite[Section~6]{HX2026} concern this already projective map:
relative vanishing transports the non-klt closed subspace;
$\alpha$ descends to a nef NQC class on $Z$; if the non-klt
subspace dominates $Z$, its contraction descends by the equality
$\mathcal O_Z=f_*\mathcal O_{\operatorname{Nklt}}$; otherwise
the projective canonical bundle formula and the multiplier-ideal
calculation identify its image with the non-klt subspace on the
base. The assertion $\mathsf C_{d-1}$ then applies to $Z$.
This proves $\mathsf C_d$ without invoking a generalized termination theorem.

\smallskip\noindent
\emph{Order of the induction.}
The assertions in dimensions zero and one are immediate from
the degree of a class on a compact curve (and the ordinary curve
minimal model program). For $d=2,3$, prove successively
\[
 \mathsf C_{d,\mathrm{big}},\quad
 \mathsf B_d,\quad \mathsf G_d,\quad\mathsf C_d.
\]
At $d=4$ prove only $\mathsf C_{4,\mathrm{big}}$ and
$\mathsf B_4$. The ordinary rational termination invoked at the
intermediate stages has dimension at most three. This closes
the induction without using \cite[Theorems~1.5 or~5.15]{HX2026}.
\end{proof}
\subsection{Termination of ordinary programs}\label{repair:sec:reduction}

\begin{corollary}\label{repair:cor:existing}
Using Theorem~1.1 of \cite{OAI:Termination-of-generalized-log-canonical-flips-on-compact-Kahler-fourfolds-October-5-2026}, every existing sequence of ordinary
rational klt flips in the global Weil $\Q$-factorial compact K\"ahler
fourfold category terminates, provided both sides of every small
diagram are projective and the ordinary adjoints have the specified
opposite ample signs. In this first formulation, compatible actual
canonical divisor representatives are fixed on the models. The contraction
bases are required to be normal compact K\"ahler, and the morphisms to
have connected fibres.
\end{corollary}
\begin{proof}
Take the fixed nef b-divisor in \cite{OAI:Termination-of-generalized-log-canonical-flips-on-compact-Kahler-fourfolds-October-5-2026} to be zero, represented on
any projective log resolution of the initial space. Its class is nef.
Klt implies log canonical;
the boundary is rational and is transported by strict transform. All
the remaining conditions in that theorem are exactly the hypotheses
just stated. No pseudo-effectivity or scaling parameter is needed.
\end{proof}

\subsubsection{The intrinsic canonical-sheaf formulation}
The printed theorem in \cite{OAI:Termination-of-generalized-log-canonical-flips-on-compact-Kahler-fourfolds-October-5-2026} specifies actual canonical divisors.
For zero nef part, its proof also has the following intrinsic
formulation; this observation is necessary to preserve the intrinsic
formulation of Theorem~\ref{thm:finite}.

For each finite comparison diagram, choose a sufficiently divisible
integer $m$ clearing the boundary denominators and the adjoint indices
on its models, and put
\[
 \mathcal D_{T,m}
   =\bigl(\omega_T^{[m]}\otimes\cO_T(mB)\bigr)^{**}.
\]
It is an invertible sheaf. On a common resolution of two models, the
canonical identification on their common smooth open defines a
distinguished meromorphic section of
\[
 p^*\mathcal D_{T,m}\otimes(q^*\mathcal D_{T',m})^{-1}.
\]
The relative canonical Jacobians and the boundary equations give its
finite meromorphic orders along exceptional divisors. Its divisor
divided by $m$ is the actual discrepancy-change divisor. It is
independent of the chosen divisible index and is compatible with
further pullback. No global meromorphic section of $\omega_T$ has
been chosen. The integer $m$ may change from one finite diagram to
another; no uniform bound on the canonical indices of an infinite
sequence is asserted.

In the zero-nef-part proof, negativity and the support comparisons
use precisely these exceptional divisors. The local adjoint algebras
are intrinsically the direct sums of the sheaves
$f_*\mathcal D_{T,m}$, and the adjunction comparisons are comparisons
of dualizing sheaves. The relative canonical divisor occurring in the
companion's branch argument is already defined in this manner.
Its local small factorialization uses the corresponding crepant
identity of reflexive canonical sheaves, and its globalization uses
the intrinsic relative canonical divisor. The Hodge, parameter-space
and monodromy steps use constant sheaves and ordinary line bundles;
they do not require a global canonical form.

The fixed denominator in the corrected difficulty clears the finitely
many initial boundary coefficients and the fixed extracted-label
coefficients, exactly as in \cite{OAI:Termination-of-generalized-log-canonical-flips-on-compact-Kahler-fourfolds-October-5-2026}. It is separate from the
diagram-dependent adjoint indices above. In the smooth-centre
positive-surface witness calculation with zero nef part, the formula
$a^+=2-\operatorname{ord}_E(B)$ needs only the initial boundary
denominators. Neither that calculation nor the corrected difficulty
requires a uniform canonical Cartier index along the sequence.
Replacing the canonical representatives by these canonical local
comparisons therefore leaves that proof unchanged and proves the
intrinsic version of Corollary~\ref{repair:cor:existing}.

\subsubsection{The supporting contraction statement}

\begin{proposition}[Supporting contractions]\label{repair:prop:supporting-contraction}
Let $(T,B)$ be an ordinary rational klt pair on a normal globally Weil
$\Q$-factorial compact K\"ahler space of dimension $1\leq d\leq4$,
with canonical sheaf a rational line bundle. Put $D=K_T+B$. For every $D$-negative extremal
ray $R$ and every supporting class with the K\"ahler margin
\[
 \alpha\text{ nef},\qquad \NA(T)\cap\alpha^\perp=R,
 \qquad \alpha-c_1(D)\text{ K\"ahler},
\]
there is a proper surjective morphism $f:T\to Z$ with connected
fibres onto a normal compact K\"ahler space and a K\"ahler class
$\gamma$ on $Z$ such that $\alpha=f^*\gamma$.
Moreover $f$ is projective and $-D$ is $f$-ample.
\end{proposition}
\begin{proof}
Write $\kappa=\alpha-c_1(D)$ and regard its pullbacks as a nef
b-class. The generalized pair $(T,B+\overline\kappa)$ is gklt:
the nef datum descends on $T$, so its discrepancies are exactly
those of $(T,B)$. Its boundary-plus-nef class $[B]+\kappa$ is big.
The assertion $\mathsf B_d$, proved in
Proposition~\ref{repair:prop:restricted-bpf-induction}, gives $f$ and
$\gamma$. This assertion applies to gklt spaces without a strong
factoriality hypothesis, as explained in that induction by a
projective small factorialization followed by descent.
Lemma~\ref{repair:lem:polarization}, with detector $D$, makes $-D$
relatively ample. The equality of the null cone with $R$ is exactly
the required contraction property.
\end{proof}

Lemma~\ref{lem:absolute-support} constructs the supporting classes by the ordinary
cone theorem. The proposition above supplies the required contraction. It can be applied after every finite prefix of the
ordinary program.

\begin{theorem}[Finite ordinary programs]
\label{repair:thm:reduction}
Let $(X,B)$ be a rational klt pair on a normal globally Weil
$\Q$-factorial compact K\"ahler fourfold. The canonical datum may
be a specified canonical Weil divisor or the intrinsic canonical
rational line bundle. Then every maximal ordinary negative-ray
program starts on $X$ and terminates. All models remain normal,
compact K\"ahler, globally Weil $\Q$-factorial and klt; global
strong $\Q$-factoriality is preserved when imposed initially.
Every negative contraction is projective, with relative
line-bundle rank and relative Bott--Chern dimension one.

If $K_X+B$ is pseudo-effective, the endpoint has nef adjoint. The
composite extracts no divisor, discrepancies do not decrease, and
strictly increase for every prime on $X$ that is contracted. If
$K_X+B$ is not pseudo-effective, the endpoint has a projective
Mori fibre contraction to a normal compact K\"ahler space of
smaller dimension, with connected fibres and relatively
antiample adjoint. In particular, these conclusions imply Theorem~\ref{thm:finite}.
\end{theorem}
\begin{proof}
Start on the given pair. If its adjoint is not nef, choose a negative
extremal ray. Proposition~\ref{repair:prop:supporting-contraction} supplies its contraction;
Lemma~\ref{repair:lem:polarization} supplies the global relative polarization
from the ordinary rational adjoint. The remaining proof of Proposition~\ref{prop:ordinary-step} gives the rational flip when the contraction
is small, and preserves klt singularities, the global factoriality
category, canonical data and the negative-side numerical ranks.
These steps use the projective analytic ordinary flip theorem
\cite{Fujino2022} and actual line-bundle descent. In the globally Weil
case this is the argument of Proposition~\ref{prop:ordinary-step}: the rational
adjoint itself gives the global flip algebra, and the transform of
each global prime is made rationally Cartier by adjusting it by a
rational multiple of that adjoint and applying integral descent.
The same argument applied to global rank-one reflexive sheaves
preserves the strong category. Relative vanishing and cohomological
descent give the one-dimensional negative-side quotient. The
argument can be repeated after each finite prefix.

Let $c_3(T)$ be the dimension of the span in $H_6(T,\R)$ of compact
analytic three-dimensional cycle classes. It is finite. A small
transformation preserves it, while a divisorial contraction strictly
decreases it: an exceptional divisor has nonzero class detected by
the cube of a K\"ahler form, and its image has dimension at most two.
The Borel--Moore localization argument is Lemma~\ref{lem:cycle-loss}.
There are therefore only finitely many divisorial steps in a run.
An infinite run would have an infinite flip tail, contrary to
Corollary~\ref{repair:cor:existing}, or its intrinsic formulation above.
Thus the run is finite.

A maximal run ends with a nef adjoint or a Mori fibre contraction,
since each negative ray otherwise admits continuation. Across each
birational step the comparison on a common resolution has the form
\[
 p^*(K_T+B)=q^*(K_{T'}+B')+E,\qquad E\geq0
\]
with $E$ exceptional over $T'$. Positive-current pullback and
exceptional descent preserve pseudo-effectivity. A relatively
antiample adjoint on a Mori fibre space is not pseudo-effective:
restrict a putative positive current to a fibre curve through a
point where a local potential is finite, and its nonnegative degree
contradicts relative antiampleness. A nef adjoint is pseudo-effective.
These observations distinguish the two endpoint cases exactly as
in Section~\ref{sec:endpoints}. Finally, the effective discrepancy
comparisons add along the finite run, with strict increase at each
contracted original divisor. They give the stated minimal-model
discrepancy inequalities.
\end{proof}
